% Public mathematical manuscript source.
\begin{abstract}
本文研究指定 Heston 模型及其原始方差正部 Euler 格式下、具有十二个月度观察时点的有界算术平均亚式看涨价差
\(\psi(A)=(A-95)^+-(A-110)^+\)，并构造有限误差证书。共同高斯因子使非负收益余项可由加权二阶矩控制。经验证的变换运算计入投影修正、积分尾部、反演误差、复对数分支及舍入误差。在
\(\theta_*=(3,.045,.23,-.55,.045)\)、\(h=1/768\) 时，Euler 价格减连续模型价格的带符号差值位于
\[
[-0.011024692273,\;0.010642371599],
\]
连续模型下的看涨价差价格位于
\[
[6.508371733,\;6.518868974].
\]
另一初始方差参数点也具有完整证书。步长比较揭示了由收益余项主导的区间宽度平台；本文还推导出使精确非线性贡献最小的共同几何尺度。两项辅助结果分别适用于不同参数域：当
\(\xi=0\) 时，九个看跌期权和亚式价差具有共同的一阶展开公式及显式余项界；在
\(\xi\in[10^{-7},10^{-6}]\) 的受限连续先验下，后验目标分位数的位移具有认证上界。本文的贡献在于加权收益结构，以及针对指定原始离散核的完整有限误差核算，并由可执行数值包提供支撑。
\end{abstract}

\section{引言}

衍生品的数值估值涉及两个不同问题：计算模型价格，以及证明计算所用离散模型对连续模型的近似精度。对于具有固定观察时点的算术平均亚式期权，正部收益函数、随机方差和离散方差投影同时存在。具有实际用途的结论应在指定步长下给出以价格为单位的误差，并在统一误差预算中计入收益转换、变换计算、无穷尾部及舍入误差。

本文围绕这一要求组织理论。我们研究风险中性 Heston 模型及一个明确规定的数值格式：对方差 Euler 更新施加正部投影，并在对数价格更新中使用当前方差。研究对象是实际算术平均收益函数在两个模型下的期望差。方法从收益函数入手：识别所有观察价格共享的高斯随机因子，对其精确积分，再利用其余变量的结构控制误差。关键量不是某个时刻方差的负阶矩，而是
\[
E\!\left[(A-cG)^2 I^{-1/2}\right],
\]
其中 \(A\) 为算术平均，\(G\) 为几何平均，\(I\) 为首个观察区间内的积分方差。

条件化与几何参考变量在亚式期权定价中已有成熟应用 \cite{Curran1994,RogersShi1995,Thompson2002,FusaiKyriakou2016}。特别地，Fusai 与 Kyriakou 的框架覆盖 Heston 模型，并给出基于条件方差的误差上界 \cite{FusaiKyriakou2016}。Heston 模型的仿射结构提供了变换工具 \cite{Heston1993}；离散几何平均亚式期权也有递归变换方法 \cite{KimKimKimWee2016} 及公开实现 \cite{QuantLibDiscreteGeometricHeston}。

主要结果将加权收益余项与原始投影 Euler 核的完整变换误差控制结合。九十一个实载荷控制非线性项，十三个复载荷确定线性项。弱展开与后验传递则是辅助应用。

本文的逻辑结构为
\[
\begin{gathered}
\left.
\begin{array}{l}
\mathsf A:\ \text{deterministic positive integrated variance}\\
\mathsf C:\ \text{verifiable Heston/projected Euler conditions}
\end{array}
\right\}\Longrightarrow\mathsf H,\\[4pt]
\mathsf H\xRightarrow{\ \mathcal S\ }\mathsf M\Longrightarrow\mathsf B .
\end{gathered}
\tag{1.1}
\]
其中 \(\mathcal S\) 表示条件高斯平滑与严格变换误差核算；\(\mathsf M\) 包括带显式常数的收益分解及变换误差界；\(\mathsf B\) 是实际价格差的严格包络区间。第~2 节精确定义这些对象，第~3--5 节证明主要推理链。第~6 节和第~7 节分别讨论弱展开与后验传递。第~8 节给出数值证书及其金融用途，第~9 节说明适用范围。


\subsection{与条件化亚式期权界的比较}


Fusai 与 Kyriakou~\cite{FusaiKyriakou2016} 已给出 Heston 亚式期权的优化条件化界（第~2.2--3 节）、条件方差误差界（定理~4，式~(51)--(54)），以及多种观察频率下的实验（第~5.1 节，表~5）。下表列出本文研究的具体数学与计算对象。本文的主要结果是针对原始核~(2.2) 的有限误差包络；其依据是由首期积分方差加权的收益余项，以及对每项所列误差贡献的验证计算。





表~1 的页码采用文献~\cite{FusaiKyriakou2016} 所链接录用稿在仓库封面之后的内部编号，而非期刊页码。

\noindent\textbf{表~1. 与 Fusai 和 Kyriakou 条件化亚式期权界的比较。}
\begin{tabularx}{\textwidth}{@{}p{.20\textwidth}XX@{}}
\toprule
对象 & Fusai--Kyriakou（已核查文献） & 本文结果 \\
\midrule
收益近似 &
优化条件下界与条件方差误差（定理~4，第~15 页）。 &
精确线性参考分解，非负余项由加权二阶矩控制（引理~3.1）。 \\
\addlinespace
离散数值分布 &
QE 模拟基准（附录~B，第~22 页）；未报告核~(2.2) 的证书。 &
控制代数变换递归与指定原始方差正部 Euler 分布之间的投影修正，股票更新使用当前方差（附录~B）。 \\
\addlinespace
变换计算 &
傅里叶条件矩（式~(55)--(56)，第~15--16 页）及 FFT 实现（第~5 节，第~16--17 页）。 &
有限实、复载荷目录，认证求积与无穷尾部界，以及复对数分支检查（附录~A--C）。 \\
\addlinespace
有限精度运算 &
已核查文献未报告向外舍入的误差预算。 &
有理端点组合与球算术包络将舍入计入其他有限误差（附录~A--C 和~G）。 \\
\addlinespace
认证输出 &
下界近似误差（定理~4）；未报告可执行的、带符号的 Euler 价格减连续价格包络。 &
指定参数与网格下的带符号 \(p_h-p_c\) 区间（定理~4.1 和推论~4.2），可追溯至配套可执行程序包。 \\
\bottomrule
\end{tabularx}


本文研究的增量在于加权余项，以及将其与原始投影核的误差控制结合。第~4 节分析由此所得证书的宽度。

\subsection{三个不同的适用域}

所有以价格为单位的亚式期权结论均针对式 (2.3) 的算术平均看涨价差，其收益上界为 15。表~2 区分三组结果的适用域。后验例子属于接近确定性模型的微扰区域，而非主 Heston 参数点的校准稳定性结论。

\noindent\textbf{表~2. 三组认证结果的适用域。}
\begin{center}
\begin{tabularx}{\textwidth}{@{}l X@{}}
\toprule
结果 & 适用域与含义 \\
\midrule
主价格证书 & 式 (4.3) 下，\(h=1/384,1/768,1/1536\)；另一 \(v_0=.04\) 参数点（第~8.3 节）仅取 \(h=1/768\)。两者均使用 \(\xi=.23\)、十二个月度观察时点及原始投影 Euler 核。 \\
十个收益函数的弱展开 & 定理 6.1 满足式 (6.1)；数值参数族 (6.10)：\(\xi=0\)、\(\kappa=3\)、\(\bar v=.045\)、\(v_0\in[.03,.06]\)，十二个月度观察时点，\(h\le1/768\) 与观察时点对齐。 \\
后验分位数传递 & 独立均匀先验 (7.5)，\(\xi\in[10^{-7},10^{-6}]\)，附录 F 中的精确合成报价。 \\
\bottomrule
\end{tabularx}
\end{center}

\section{模型与结构条件}

\subsection{连续模型、离散核与收益函数}

设 \(W,B\) 为独立标准布朗运动，且
\(\theta=(\kappa,\bar v,\xi,\rho,v_0)\)。
假设
\(S_0,\kappa,\bar v,v_0>0\)、\(\xi\ge0\)、\(|\rho|<1\) 和 \(r\in\mathbb R\)。
离散模型使用 \(h>0\)，其中每个 \(t_i/h\) 均为正整数。连续模型为
\[
\begin{aligned}
dV_t&=\kappa(\bar v-V_t)\,dt+\xi\sqrt{V_t}\,dW_t,\\
dZ_t&=(r-V_t/2)\,dt+
\sqrt{V_t}\bigl(\rho\,dW_t+\sqrt{1-\rho^2}\,dB_t\bigr),\\
S_t&=S_0e^{Z_t},\qquad Z_0=0,\quad V_0=v_0.
\end{aligned}
\tag{2.1}
\]
在与观察时点对齐、步长为 \(h\) 的网格上，定义原始离散核 \(Q_h\) 为
\[
\begin{aligned}
V_{j+1}^h&=\left[(1-\kappa h)V_j^h+\kappa\bar v h+
\xi\sqrt{hV_j^h}\,G_j\right]^+,\\
Z_{j+1}^h&=Z_j^h+(r-V_j^h/2)h+
\sqrt{hV_j^h}\bigl(\rho G_j+\sqrt{1-\rho^2}H_j\bigr),
\end{aligned}
\tag{2.2}
\]
其中 \(G_j,H_j\) 为独立标准正态变量，同一步的两个更新共享 \(G_j\)，初始条件一致。特别地，价格更新使用 \(V_j^h\)。

固定 \(0<t_1<\cdots<t_n=T\)，并令
\[
A=\frac1n\sum_{i=1}^nS_{t_i},\qquad
G=\left(\prod_{i=1}^nS_{t_i}\right)^{1/n},\qquad
\psi(x)=(x-K_1)^+-(x-K_2)^+,
\tag{2.3}
\]
其中 \(0<K_1<K_2\)。记
\[
p_c=e^{-rT}E_P\psi(A),\qquad
p_h=e^{-rT}E_{Q_h}\psi(A),\qquad e_h=p_h-p_c.
\tag{2.4}
\]
当步长明确时，简记为 \(Q=Q_h\)。全文采用这一误差符号约定。数值例子使用 \(n=12,t_i=i/12,T=1\)，且
\[
S_0=100,\quad r=.01,\quad K_1=95,\quad K_2=110,\quad h=1/768.
\tag{2.5}
\]
校准收益函数为九个欧式看跌期权，到期时间为 \(1/4,1/2,1\)，执行价为 \(90,100,110\)。

\subsection{共同高斯结构 \(\mathsf H_s\)}

称定价分布 \(\nu\) 满足 \(\mathsf H_s(c)\)，如果存在子 \(\sigma\) 代数 \(\mathcal G\)、正的 \(\mathcal G\) 可测随机变量 \(a_1,\ldots,a_n,\sigma\)，以及随机变量 \(U\)，使得
\[
U\mid\mathcal G\sim N(0,\sigma^2),\qquad S_{t_i}=e^Ua_i.
\tag{2.6}
\]
令 \(a=n^{-1}\sum_i a_i\) 和 \(g=(\prod_i a_i)^{1/n}\)。我们要求
\[
E_\nu(A+G)<\infty,\qquad
D_\nu(c):=E_\nu\!\left[\frac{(A-cG)^2}{\sigma}\right]<\infty.
\tag{2.7}
\]
其中 \(c>0\) 为预先选定的常数。条件~(2.6) 描述模型的随机结构，而条件 (2.7) 是可通过矩或 Laplace 变换核查的可积性条件。两者均未预先包含有待证明的价格差结论。

\subsection{有限验证条件 \(\mathsf H_v\)}

结构条件还须与有限计算建立联系。本文所用条件 \(\mathsf H_v\) 包含下列输入：
\begin{enumerate}
\item \(D_\nu(c)\) 的可计算上界，由有限 Laplace 节点上的包络、可证明的求积包络及无穷尾部界共同给出。
\item 第~3 节所需有限变换载荷集的向外包络。对离散核，每个载荷还具有真核与其代数递归之间差值的可计算上界。
\item 傅里叶反演的无穷频率尾部与周期化误差的完整上界，以及连续复变换的分支控制。
\end{enumerate}
这些输入均由具体变换或矩定义；以下定理建立它们与目标价格的联系。记
\(\mathsf H=\mathsf H_s+\mathsf H_v\)。
有限目录、精度、尾部界常数及核误差公式详见附录~A--C。

\subsection{中间性质与目标结论}

性质 \(\mathsf M\) 包含两部分：将实际收益精确分解为线性项与余项，以及每个余项和每个线性变换误差的带符号包络。目标结论 \(\mathsf B\) 是明确的有限区间 \([\underline e,\overline e]\)，使得
\[
e_h\in[\underline e,\overline e].
\tag{2.8}
\]
每个端点均由有限次有理运算和超越函数的经验证包络得到。

\section{核心论证：从共同高斯因子到加权余项}

\subsection{条件平滑引理}

\textbf{引理 3.1.} 假设 \(\mathsf H_s(c)\) 成立。对 \(K>0\)，定义
\[
R_{K,\nu}=E_\nu\!\left[(A-K)^+-(cG-K)^+
-(A-cG)\mathbf1_{\{cG>K\}}\right].
\tag{3.1}
\]
则
\[
0\le R_{K,\nu}\le \frac{D_\nu(c)}{2K\sqrt{2\pi}}.
\tag{3.2}
\]

\textit{证明。} 正部函数的凸性给出逐路径非负性。在给定 \(\mathcal G\) 的条件下，令
\[
F_K(x)=E[(xe^U-K)^+\mid\mathcal G],\qquad x>0.
\]
由于 \(\sigma>0\)，该函数二次可微，且
\[
F_K''(x)=\frac{K}{\sigma x^2}
\varphi\!\left(\frac{\log(x/K)}{\sigma}\right)
\le\frac{e^{2\sigma^2}}{K\sigma\sqrt{2\pi}},
\tag{3.3}
\]
其中 \(\varphi\) 为标准正态密度。最后一个上界由在
\(\log(x/K)=-2\sigma^2\)
处取最大值得到。

在正数 \(a,cg\) 之间应用带积分余项的 Taylor 公式，得到
\[
0\le F_K(a)-F_K(cg)-(a-cg)F_K'(cg)
\le \frac{e^{2\sigma^2}(a-cg)^2}{2K\sigma\sqrt{2\pi}}.
\]
条件导数满足
\[
(a-cg)F_K'(cg)
=E[(A-cG)\mathbf1_{\{cG>K\}}\mid\mathcal G].
\]
由于
\[
E[(A-cG)^2\mid\mathcal G]=e^{2\sigma^2}(a-cg)^2,
\]
取期望即得式 (3.2)。看涨项与线性项因 \(E(A+G)<\infty\) 而可积。条件二阶矩恒等式按非负条件期望理解，加权余项由 \(D_\nu(c)<\infty\) 控制。\hfill\(\square\)

关键在于两个因子 \(e^{2\sigma^2}\) 精确匹配。因此，条件平滑的曲率贡献可直接表示为真模型下的加权二阶矩。残差 \(A-cG\) 可为正或负。

\subsection{收益分解}

定义
\[
\mathcal L=\psi(cG)+(A-cG)\mathbf1_{\{K_1<cG\le K_2\}}.
\tag{3.4}
\]
条件正态方差严格为正，因此 \(G\) 在任何正阈值处均无原子。由式 (3.1)，
\[
E_\nu\psi(A)=E_\nu\mathcal L+R_{K_1,\nu}-R_{K_2,\nu}.
\tag{3.5}
\]
这一恒等式保留原始算术平均收益作为研究对象，同时给出可用少量变换计算的线性部分。

令 \(Y=\log(G/S_0)\) 和 \(y_K=\log(K/(cS_0))\)，并定义两个正测度的分布函数为
\[
F_{0,\nu}(y)=\nu(Y\le y),\qquad
F_{A,\nu}(y)=E_\nu[A\mathbf1_{\{Y\le y\}}].
\]
在各区间展开式 (3.4)，得到
\[
E_\nu\mathcal L=(K_2-K_1)
+K_1F_{0,\nu}(y_{K_1})-K_2F_{0,\nu}(y_{K_2})
-F_{A,\nu}(y_{K_1})+F_{A,\nu}(y_{K_2}).
\tag{3.6}
\]
相应变换为
\[
\widehat\mu_0(z)=E_\nu e^{zY},\qquad
\widehat\mu_A(z)=\frac{S_0}{n}\sum_{i=1}^n
E_\nu e^{Z_{t_i}+zY}.
\tag{3.7}
\]
因此，当 \(n=12\) 时，线性部分需要十三个载荷。

\section{主定理与完整价格证书}

\textbf{定理 4.1（两个模型价格差的带符号包络）。}
假设 \(P,Q_h\) 分别满足 \(\mathsf H_s(c)\)，且可计算常数满足
\(D_P(c)\le d_P,D_Q(c)\le d_Q\)。若 \(\mathsf H_v\) 给出
\[
e^{-rT}(E_Q\mathcal L-E_P\mathcal L)\in[\ell,u],
\tag{4.1}
\]
则令 \(\gamma=e^{-rT}/(2\sqrt{2\pi})\)，有
\[
\boxed{
e_h\in
\left[
\ell-\gamma\!\left(\frac{d_Q}{K_2}+\frac{d_P}{K_1}\right),
\quad
u+\gamma\!\left(\frac{d_Q}{K_1}+\frac{d_P}{K_2}\right)
\right].}
\tag{4.2}
\]

\textit{证明。} 分别对两个分布应用式 (3.5) 并相减。下端点贡献为
\(-R_{K_2,Q}-R_{K_1,P}\)，上端点贡献为
\(R_{K_1,Q}+R_{K_2,P}\)。引理~3.1 与式 (4.1) 给出结论。\hfill\(\square\)

主定理的最后一步仅需组合区间。其数学实质在于引理~3.1 与 \(\mathsf H_v\) 的严格验证。所得区间通常不关于零对称，计算中保留其中心。


\textbf{推论 4.2（指定 Heston 参数点）。}
对式 (2.1)--(2.5)，取
\[
\theta_*=(3,\;9/200,\;23/100,\;-11/20,\;9/200),\qquad
c=\frac{1254433}{1250000}.
\tag{4.3}
\]
则
\[
e_h\in[-.011024692273,\;.010642371599],\qquad |e_h|<.011025<.025.
\tag{4.4}
\]
此外，
\[
\begin{aligned}
p_c&\in[6.508371733,\;6.518868974],\\
p_h&\in[6.507844281,\;6.519014106].
\end{aligned}
\tag{4.5}
\]

\textit{证明。} 第~5 节建立共同高斯结构。附录~A--C 列明经验证运算的全部要求，计算得到
\[
W_P:=E_P[(A-cG)^2I^{-1/2}]\le2.172471592802,\qquad
W_Q\le2.242189951600.
\tag{4.6}
\]
其中 \(d_\nu=W_\nu/\sqrt{1-\rho^2}\)。有限线性差值的中心约为
\(-.000202992309237006\)，其完整误差半径至多为
\(.000594238060\)。根据式 (4.2) 向外组合精确有理端点，得到式 (4.4)。分别对两个分布的式 (3.5) 应用相应误差贡献，得到式 (4.5)。源代码、精确端点，以及第二套实现对整个目录的重算，共同构成附录所规定的可执行证据。\hfill\(\square\)

\subsection{加权余项的优化及其宽度局限}

共同尺度 \(c\) 影响非线性转换界。除结构条件外，进一步假设
\[
m_{AA,\nu}=E_\nu[A^2/\sigma_\nu],\qquad
m_{AG,\nu}=E_\nu[AG/\sigma_\nu],\qquad
m_{GG,\nu}=E_\nu[G^2/\sigma_\nu]
\]
对 \(\nu=P,Q\) 均有限。令
\[
\begin{aligned}
a&=m_{AA,P}+m_{AA,Q},\qquad b=m_{AG,P}+m_{AG,Q},\\
d&=m_{GG,P}+m_{GG,Q}>0.
\end{aligned}
\]
采用精确矩输入时，式 (4.2) 中非线性部分对宽度的贡献为
\[
\gamma\left(K_1^{-1}+K_2^{-1}\right)J(c),\qquad
J(c)=a-2bc+dc^2.
\]
因此，在 \(c>0\) 上的唯一极小点为
\[
c_* = \frac{b}{d},\qquad
J(c)=J(c_*)+d(c-c_*)^2.                         \tag{4.7}
\]
折现因子与执行价因子均为正常数，故不改变极小点。该结论针对精确非线性宽度；求积松弛量与线性变换误差容限可能依赖于 \(c\)，因此并不说明 \(c_*\) 使最终计算包络最小。

在 Heston 构造中，\(\sigma_\nu=\sqrt{1-\rho^2}\sqrt{I_\nu}\)。
共同因子 \((1-\rho^2)^{-1/2}\) 在式 (4.7) 中消去。数值上更合适的表达利用算术平均与几何平均不等式：
\[
H_\nu=E_\nu[G^2I_\nu^{-1/2}],\qquad
E_\nu^+=E_\nu[(A-G)GI_\nu^{-1/2}]\ge0,
\qquad
c_* =1+\frac{E_P^++E_Q^+}{H_P+H_Q}.             \tag{4.8}
\]
两个矩均使用 91 个载荷组合目录中已有的组合。例如，Laplace 被积函数
\(E_\nu[(A-G)G\exp(-x^2I_\nu)]\) 非负且单调不增。
因此，可按式 (A.5) 的同一 Darboux 构造给出包络：在积分前，逐节点计算 \(AG\) 与 \(G^2\) 载荷组合之差。其连续尾部由正的
\(AG\) 尾部控制。在 \(S_0=100\) 时，每个正原始矩目录的系数和为
\(10000\)；\((A-G)G\) 的节点投影容限为
\(20000\delta_Q\)。这些做法在节点层面保留抵消，避免将两个宽积分矩区间相减。

若经验证的正端点给出
\(E_P^++E_Q^+\in[e_-,e_+]\) 和 \(H_P+H_Q\in[d_-,d_+]\)，且
\(d_->0\)，则
\[
c_*\in\left[1+\frac{e_-}{d_+},\;1+\frac{e_+}{d_-}\right]=[c_-,c_+].
\]
选择有理中点 \(\widehat c=(c_-+c_+)/2\) 得到
\[
0\le J(\widehat c)-J(c_*)
\le \frac{d_+(c_+-c_-)^2}{4\sqrt{1-\rho^2}},    \tag{4.9}
\]
其中 \(d_+\) 包络两个以 \(I^{-1/2}\) 加权的
\(G^2\) 矩之和，如式 (4.8) 所示。选定的有理尺度仍须代入原始平方项和线性项计算；单独的比值包络并不构成新的价格证书。

优化 \(c\) 无法恢复两个模型之间的抵消。即使
\(Q=P\)、精确线性差为零，且 \(d_P=d_Q=D_P(c)\) 精确，式 (4.2) 的宽度仍为
\[
2\gamma\left(K_1^{-1}+K_2^{-1}\right)D_P(c).
\]
除非 \(A=cG\) 几乎处处成立，否则该量为正。因此，式 (4.2) 是固定步长包络，本身并不能证明其宽度随
\(h\) 减小而趋零。


对于计算输入，式 (4.2) 的实际宽度为
\[
(u-\ell)+\gamma(d_P+d_Q)(K_1^{-1}+K_2^{-1}).
\tag{4.10}
\]
原非线性区间宽度约为 .02047859，线性宽度约为 .00118848，总宽度约为 .02166706。因此，收益转换约占百分之 94.5。提高精度或缩小已经很小的投影贡献，对最终宽度的作用有限。
附录 I 给出直接控制余项差的充分耦合不等式。其所需额外矩尚未在主参数点计算，因此推论 4.2 保留分别控制两个分布余项的做法。在附录~I 的确定性例子中，非线性宽度由约 .01538363715 降至 .00008815020，约为分别控制余项所得宽度的百分之 0.573。两种计算采用相同条件化与共同几何尺度。该比较仅针对这一例子的非线性部分。

% INSERT_SCALE_BEGIN
在 \(h=1/768\) 处的正矩计算给出
\[
c_*\in[1.003427394603466100946724,\;1.003669831096000915452372].
\tag{4.11}
\]
将中点舍入至 \(10^{-7}\) 网格，得到有理数
\(5017743/5000000=1.0035486\)。这不是式 (4.11) 有理端点的精确中点。式 (4.9) 适用于精确中点；对舍入后的选择，距离容限还须计入 \(5\cdot10^{-8}\)。

\noindent\textbf{表~3. 共同几何尺度对非线性宽度的影响。}
\begin{center}
\begin{tabularx}{\textwidth}{@{}l X@{}}
\toprule
共同尺度 & \(h=1/768\) 时的非线性宽度上界 \\
\midrule
1 & .040551110717 \\
1.0035464（原尺度） & .020478587752 \\
1.0035486（舍入中点） & .020478592980 \\
\bottomrule
\end{tabularx}
\end{center}

原尺度已接近极小点。以中点为依据的候选值使计算上界略有增加，因此保留原有理尺度。该比较计入新矩求积、投影和尾部容限，并不声称精确目标或完整价格证书已经过全局优化。

% INSERT_SCALE_END

\section{经典条件的恢复与所覆盖的其他模型}

\subsection{确定性正积分方差：\(\mathsf A\Rightarrow\mathsf H\)}

定义 \(\mathsf A\) 如下：模型的方差系数本身为确定性函数 \(v(t)\ge0\)，股票价格服从相应的线性几何扩散，且
\[
0<I_1=\int_0^{t_1}v(t)\,dt,\qquad \int_0^T v(t)\,dt<\infty.
\tag{5.1}
\]
可计算性结论还要求所有观察区间的积分方差均作为有效可计算输入提供。

取首区间的股票布朗积分为 \(U\)。此时 \(\sigma^2=I_1\) 为正常数；其余对数增量与确定性漂移均可纳入 \(\mathcal G\)。每个观察价格都具有式 (2.6) 的共同因子。联合对数正态分布的有限正矩给出
\[
D_\nu(c)=I_1^{-1/2}E_\nu(A-cG)^2<\infty.
\]
二阶矩可由下式直接计算：
\[
E(A-cG)^2=\frac1{n^2}\sum_{i,j}ES_iS_j
-\frac{2c}{n}\sum_iES_iG+c^2EG^2
\tag{5.2}
\]
其中每项均为高斯指数矩。相应有限变换与高斯尾部也具有显式界。因此，对具有有效可计算输入的单个分布，\(\mathsf A\Rightarrow\mathsf H\)；对两个这样的分布，定理~4.1 进一步给出 \(\mathsf B\)。

这表明经典对数正态条件化模型属于本框架。识别共同结构后，同一收益层面的论证即可用于不同方差机制；具体常数与计算目录由模型本身提供。

\subsection{Heston 与实际投影核：\(\mathsf C\Rightarrow\mathsf H\)}

条件 \(\mathsf C\) 要求 \(\kappa,\bar v,\xi,v_0>0\)、\(|\rho|<1\)，采用模型 (2.1) 或 (2.2)，观察时点与离散网格对齐，并满足附录~A--C 所列有限变换、矩增长、投影及无穷尾部不等式。本节的实际验证使用 \(S_0,r,T,h\)（见式 (2.5)）和式 (4.3) 的参数；下述共同高斯构造在刚才列明的结构参数假设下有效。

条件化于完整方差驱动过程 \(W\)，以及随后独立股票驱动过程的增量 \(B_t-B_{t_1}\)、\(t\ge t_1\)。条件化的 σ 代数不包括未来布朗运动水平 \(B_t\)，因为后者包含首区间终点的值。令
\[
I_P=\int_0^{t_1}V_t\,dt,\qquad
I_Q=h\sum_{j=0}^{t_1/h-1}V_j^h.
\]
首区间内 \(B\) 的贡献由所有观察时点共享。因此
\[
U\mid\mathcal G\sim N(0,(1-\rho^2)I_\nu),\qquad S_{t_i}=e^Ua_i.
\tag{5.3}
\]
对连续模型，\(v_0>0\) 与方差路径连续性蕴含 \(I_P>0\) 几乎必然成立。对离散模型，有精确下界
\[
I_Q\ge hv_0>0.
\tag{5.4}
\]
附录中的变换与尾部界验证 \(W_\nu<\infty\) 并给出上界，指数矩增长界则保证所需价格项的可积性。因此 \(\mathsf C\Rightarrow\mathsf H\)。在式 (4.3) 的参数下，所有这些不等式与有限节点检查均成立。

\subsection{所覆盖模型类的严格扩展}

在以下概率公式中，\(\Phi\) 表示标准正态累积分布函数。

令 \(X_*\) 为式 (4.3) 的连续与离散模型对。对连续模型，
\[
\operatorname{Var}(V_t)=\xi^2\int_0^t e^{-2\kappa(t-s)}EV_s\,ds>0
\quad(t>0),
\tag{5.5}
\]
因此，其方差并非确定性函数。离散模型的第一步也满足
\[
Q_h(V_1^h=0)=
\Phi\!\left(
-\frac{(1-\kappa h)v_0+\kappa\bar v h}{\xi\sqrt{hv_0}}
\right)\in(0,1).
\tag{5.6}
\]
同时，式 (5.4) 仍成立。因此，\(X_*\) 满足 \(\mathsf C\)，但不满足 \(\mathsf A\)，主定理给出式 (4.4)。新增模型对同时具有随机方差和投影产生的边界原子，而收益层面的共同高斯结构仍然成立。


\section{共同密度扰动与十个收益函数的弱展开}

弱误差展开与分布误差分析的背景可参见
\cite{TalayTubaro1990,MickelNeuenkirch2023,BallyTalay1996}。
本节直接利用高斯密度的可微性。令 \(\xi=0\)，保留原始更新 (2.2)，并假设
\[
\kappa>0,\quad \bar v,v_0>0,\quad
t_j=j/12,\quad h=1/(12m)\le h_0<1/\kappa,\quad m\in\mathbb N.
\tag{6.1}
\]
设 \(D=v_0-\bar v\)。连续与离散累计方差为
\[
I(t)=\bar vt+\frac D\kappa(1-e^{-\kappa t}),\qquad
I_h(t)=\bar vt+\frac D\kappa[1-(1-\kappa h)^{t/h}].
\tag{6.2}
\]
对数增量向量 \(Y\) 的十二个独立分量服从
\[
N(r/12-s_j/2,s_j),\qquad s_j=I(t_j)-I(t_{j-1}),
\tag{6.3}
\]
离散分布则将 \(s_j\) 替换为 \(s_{h,j}\)。
二者均满足下界
\(s_j,s_{h,j}\ge u_*:=\min(v_0,\bar v)/12>0\)。

定义
\[
a(t)=\frac{D\kappa t e^{-\kappa t}}2,
\quad
B_I(t)=|D|e^{-\kappa t}
\left[
\frac{\kappa^2t}{3(1-\kappa h_0)}
+\frac{\kappa^3t^2}{8(1-\kappa h_0)^2}
\right].
\tag{6.4}
\]
令 \(d_j=a(t_j)-a(t_{j-1})\)，且
\(b_j=B_I(t_j)+B_I(t_{j-1})\)。附录~D 证明
\[
|s_{h,j}-s_j-hd_j|\le h^2b_j.
\tag{6.5}
\]

\textbf{定理 6.1（有界收益函数的共同一阶展开）。}
在式 (6.1) 下，设 \(g(Y)\) 为任意有界 Borel 收益函数，其值域包含于
\([l,l+L]\)，并在时刻 \(\tau\) 支付。则
\[
p_h(g)-p_c(g)=h\beta_g+R_h(g),
\tag{6.6}
\]
\[
\begin{aligned}
\beta_g&=e^{-r\tau}E_s\!\left[g(Y)\sum_jd_j\mathcal S_j(Y)\right],\\
|R_h(g)|&\le e^{-r\tau}\frac L2\,C_2h^2,
\end{aligned}
\tag{6.7}
\]
其中
\[
\mathcal S_j=\frac{G_j^2-1}{2s_j}-\frac{G_j}{2\sqrt{s_j}},
\quad
J(s)=\frac1{2s^2}+\frac1{4s},\quad
H_2(s)=\frac2{s^2}+\frac1{2s},
\]
\[
C_2=
\sqrt{\sum_jb_j^2J(u_*)}
+\frac12\sum_j(|d_j|+h_0b_j)^2H_2(u_*).
\tag{6.8}
\]
这里 \(G_j\) 表示标准化独立正态增量。

\textit{证明。}
高斯密度在 \(L^1\) 中关于方差参数二次可微。一阶方向导数范数由相应 Fisher 信息二次型的平方根控制，二阶方向导数满足
\[
\|Dp_s[z]\|_1\le\sqrt{\sum_jz_j^2J(s_j)},\qquad
\|D^2p_s[z,z]\|_1\le\sum_jz_j^2H_2(s_j).
\]
将式 (6.5) 代入密度的 \(L^1\) Taylor 公式，得到
\(\|p_{s_h}-p_s-hDp_s[d]\|_1\le C_2h^2\)。
密度余项的积分为零，故可将收益函数中心化为
\(g-l-L/2\)，从而得到式 (6.7)。
附录~D 给出导数计算与基于控制收敛定理的完整论证。
\hfill\(\square\)

该定理对全部十个收益函数采用同一主导密度扰动。对于看跌期权，
\[
\beta_{K,\tau}
=a(\tau)\frac{S_0\varphi(d_1)}{2\sqrt{I(\tau)}},
\quad
d_1=\frac{\log(S_0/K)+r\tau+I(\tau)/2}{\sqrt{I(\tau)}}.
\tag{6.9}
\]
当 \(v_0\ne\bar v\) 时，九个看跌期权分量均非零，且与
\(v_0-\bar v\) 同号。

在参数族
\[
\kappa=3,\quad\bar v=.045,\quad v_0\in[.03,.06],\quad h_0=1/768,
\tag{6.10}
\]
上，原始亚式收益的二阶余项具有统一上界
\(.000165217\)。另一个有限步长密度比较还给出
\[
|p_h(\psi)-p_c(\psi)|\le .004990924.
\tag{6.11}
\]
亚式主导系数具有式 (6.7) 的精确积分表示与有效绝对值界。该共同展开通过同一可积密度扰动处理多个非光滑收益函数，并显式保留每个余项常数。

% INSERT_BETA_BEGIN
\subsection{亚式系数界与数值示例}

式 (6.7) 的得分函数界与一阶导数范数给出
\[
|\beta_{\rm Asian}|
\le {15e^{-rT}\over2}\sqrt{\sum_{j=1}^{12}d_j^2J(s_j)}.
\tag{6.12}
\]
表~4 给出数值参数族 (6.10) 的逐点系数界与余项界。后者由式 (6.5)--(6.8) 结合各点的增量方差下界得到。这提供了共同展开公式与显式余项界；带符号的十二维亚式系数积分尚未计算，系数包络均包含零。

\noindent\textbf{表~4. 亚式系数的解析包络与余项界。}
\begin{center}
\begin{tabularx}{\textwidth}{@{}l X X@{}}
\toprule
初始方差 & 认证系数包络 & \(|e_h-h\beta|\) 的上界，\(h=1/768\) \\
\midrule
0.03 & \([-3.441491358,3.441491358]\) & 0.000121667773 \\
0.04 & \([-0.928189505,0.928189505]\) & 0.000036646884 \\
0.05 & \([-0.785853082,0.785853082]\) & 0.000034532369 \\
0.06 & \([-2.057270822,2.057270822]\) & 0.000099684585 \\
\bottomrule
\end{tabularx}
\end{center}

Monte Carlo 示例在每个点使用 2,097,152 个高斯增量向量，表~5 的四个步长共享样本。按初始方差递增顺序，系数估计为 -0.041015（标准误 0.002079）、-0.011364（标准误 0.000566）、0.008329（标准误 0.000483）和 0.019917（标准误 0.001274）。这些是抽样估计，与表~4 的解析包络不同。

\noindent\textbf{表~5. 缩放后亚式展开残差的抽样估计。}
\begin{center}
\begin{tabularx}{\textwidth}{@{}l X X X X@{}}
\toprule
初始方差 & \(h=1/192\) & \(1/384\) & \(1/768\) & \(1/1536\) \\
\midrule
0.03 & -0.063243 & -0.062944 & -0.062796 & -0.062722 \\
0.04 & -0.017226 & -0.017145 & -0.017105 & -0.017085 \\
0.05 & 0.011698 & 0.011646 & 0.011620 & 0.011607 \\
0.06 & 0.027912 & 0.027785 & 0.027722 & 0.027690 \\
\bottomrule
\end{tabularx}
\end{center}

表中数值通过高斯密度比计算估计 \((e_h-h\beta_{\rm Asian})/h^2\)。在 \(h=1/768\) 时，估计标准误分别为 .003733、.001000、.000863 和 .002319；同一点各步长的估计相互相关。数值尺度的稳定性与二阶展开一致。\(h=1/192\) 和 \(1/384\) 两列为诊断结果，超出固定 \(h_0=1/768\) 范围；该范围用于式 (6.10) 和表~4 的数值常数。

% INSERT_BETA_END

\section{从参考参数族到随机波动率后验分位数}

\subsection{波动率的波动率取小正值时的价格传递}

假设确定性参考方差及其 Euler 值位于 \([m,M]\)，其中 \(0<m\le M<\infty\)、\(\kappa\ge\kappa_->0\)，并假设
\[
\kappa h\le1,\qquad
0\le\xi\le\sqrt{\kappa_-M/2}.
\tag{7.1}
\]
定义
\[
C_P=\sqrt{M/(2\kappa_-)},\qquad C_Q=\sqrt{2M/\kappa_-}.
\]
连续模型的 It\^o 等距性与离散正部投影的 1-Lipschitz 性分别给出
\[
E(V_t^\xi-v(t))^2\le\xi^2C_P^2,\qquad
E(V_j^{\xi,h}-v_j)^2\le\xi^2C_Q^2.
\tag{7.2}
\]
将确定性参考下界 \(m\) 与 Doob 不等式结合，得到
\[
E\max|Z_\ell^\xi-Z_\ell^0|
\le\xi C_\ell\left(\frac T2+2\sqrt{T/m}\right),\quad \ell=P,Q.
\tag{7.3}
\]
原始看跌收益关于对数股票价格为 \(K\)-Lipschitz；原始价差收益关于观察对数价格向量的最大范数为
\(K_2=110\)-Lipschitz。
因此，式 (7.3) 给出所有真实校准价格及真实目标价格到其参考参数族对应价格的统一半径。附录~E 给出归纳证明，并保留投影对均值的影响。

例如，在具有正五维体积的参数盒上，
\[
\begin{gathered}
\kappa\in[2.999,3.001],\quad\bar v\in[.04499,.04501],\\
v_0\in[.03001,.05999],\quad \xi\in[10^{-5},4\cdot10^{-5}],
\quad\rho\in[-.8,-.3],
\end{gathered}
\tag{7.4}
\]
参数坐标顺序为 \((\kappa,\bar v,\xi,\rho,v_0)\)。有限验证给出原始亚式收益的绝对偏差界 \(.020948\)，以及每个看跌期权的绝对偏差界 \(.016702\)。

\subsection{完整联合后验}

前向数值误差对似然及其归一化常数的影响是贝叶斯近似中的基本问题~\cite{CotterDashtiStuart2010}。
本节令 \(\kappa=3,\bar v=.045\)，并设
\[
u=v_0\sim U[.03,.06],\quad
\xi\sim U[10^{-7},10^{-6}],\quad
\rho\sim U[-.8,-.3]
\tag{7.5}
\]
相互独立。
记 \(\zeta=(\xi,\rho)\)，以 \(\nu\) 表示其先验概率测度。在本节后验计算及附录~F 中，\(m\in\{P,Q\}\) 用于区分定价模型；\(\nu\) 仍表示辅助参数先验。
对模型 \(m=P,Q\)，令 \(c_m(u,\zeta)\) 为真实校准向量，\(J_m(u,\zeta)\) 为真实亚式价格。给定合成报价 \(y\) 与正定协方差矩阵 \(\Sigma_d\)，定义
\[
w_m(u,\zeta)=
\exp\!\left[-\frac12(c_m-y)^T\Sigma_d^{-1}(c_m-y)\right],
\quad \bar w_m(u)=\int w_m(u,\zeta)\nu(d\zeta),
\]
\[
Z_m=\frac1{.03}\int_{.03}^{.06}\bar w_m(u)\,du,\qquad
\mu_m(du,d\zeta)=\frac{w_m(u,\zeta)}{.03Z_m}\,du\,\nu(d\zeta).
\tag{7.6}
\]
每个模型使用各自的校准价格、归一化常数与目标价格。

\subsection{适用于所有分位水平的传递定理}

\textbf{定理 7.1.}
假设 \(u\) 的两个边缘后验存在耦合 \((U_P,U_Q)\)，使得 \(|U_P-U_Q|\le\delta_u\) 几乎必然成立。
假设参考目标满足
\[
|J_P^0(u)-J_P^0(v)|\le L_0|u-v|,
\quad
|J_Q^0(u)-J_P^0(u)|\le\varepsilon_0,
\]
且真实目标满足
\(|J_m(u,\zeta)-J_m^0(u)|\le e_{m,J}\)。
则真实目标在各自完整联合后验下的左广义分位数满足
\[
\boxed{\quad
|q_Q(p)-q_P(p)|
\le \varepsilon_0+e_{P,J}+e_{Q,J}+L_0\delta_u
\quad(0<p<1).\quad}
\tag{7.7}
\]

\textit{证明。}
对每个模型，定义条件概率核
\[
K_m(u,d\zeta)=\frac{w_m(u,\zeta)}{\bar w_m(u)}\nu(d\zeta).
\]
令 \(\zeta_P\) 和 \(\zeta_Q\) 分别按 \(K_P(U_P,\cdot)\) 和 \(K_Q(U_Q,\cdot)\) 抽取，以扩展上述边缘耦合。
Fubini 定理表明，所得完整边缘分布恰为 \(\mu_P\) 与 \(\mu_Q\)。三角不等式给出
\[
|J_Q(U_Q,\zeta_Q)-J_P(U_P,\zeta_P)|
\le\varepsilon_0+e_{P,J}+e_{Q,J}+L_0\delta_u=:E.
\]
因此，\(F_P(z-E)\le F_Q(z)\le F_P(z+E)\)。
取左广义逆即得式 (7.7)。论证同样适用于含原子的目标分布。
\hfill\(\square\)

\subsection{有效输入与数值结论}

将 \([.03,.06]\) 分成 4096 个等宽区间。在每个区间及整个 \(\zeta\) 域上，用式 (7.3) 与参考 Black 价格包络真实似然。对归一化先验 \(\nu\) 积分后，相同上界仍成立。在保留共同分母的同时构造边缘累积分布函数包络；由共同均匀变量驱动的逆累积分布函数耦合给出 \(\delta_u\)。所有常数与有限配对算法详见附录~F。

单报价情形使用 \(T=1,K=100\) 的看跌期权，九报价情形使用完整校准向量。两者均取
\(\Sigma_d=.75^2(.75I_d+.25\mathbf1\mathbf1^T)\)。
精确报价列于附录~F。结果为
\[
L_0\le213.792201829985,\quad
\varepsilon_0\le.004493504474,\quad
e_{P,J}+e_{Q,J}\le.000393595478,
\]
\[
\begin{array}{c|c|c}
\text{Case}&\delta_u&\sup_{0<p<1}|q_Q(p)-q_P(p)|\\ \hline
\text{Single quote}&3/204800&\le.008018821658\\
\text{Nine quotes}&3/81920&\le.012716404217
\end{array}
\tag{7.8}
\]
有限分割包络连续后验，并保留每个模型的完整归一化质量。
\section{数值证书与敏感性}

\subsection{从公式到货币单位下的认证界}

表~6 汇总原始参数点的价格差。所列包络和上界均向外舍入。
精确端点、数值构造与复算说明见附录~G。

\noindent\textbf{表~6. 主 Heston 价格证书的误差预算。}
\begin{center}
\begin{tabularx}{\textwidth}{@{}X r@{}}
\toprule
量 & 经验证的上界或包络 \\
\midrule
非线性转换区间 & \([-.010227461904,.010251125849]\) \\
Q 模型每个线性载荷的投影误差 & \(.000000012711979137\) \\
Q 模型全部线性载荷的投影贡献 & \(.000004310413\) \\
连续模型的截断频率尾部 & \(.000014095712\) \\
离散模型的截断频率尾部 & \(.000184373301\) \\
单个模型的周期化误差 & \(.000195729318\) \\
两模型线性差的总误差界 & \(.000594238060\) \\
\textbf{真实算术亚式期权偏差的绝对界} & \(\boldsymbol{.011025}\) \\
\bottomrule
\end{tabularx}
\end{center}

% INSERT_STEPS_BEGIN
所列宽度由精确误差账本计算后向外舍入。
直接相减显示的端点，末位小数可能不同。
对 \([\ell,u]\)，宽度为 \(u-\ell\)，中点半径为
\((u-\ell)/2\)，绝对偏差界为
\(\max(|\ell|,|u|)\)。

\subsection{步长网格证书与宽度平台}

表~7 保留主 Heston 参数、收益函数、观察日期和变换目录，
仅改变Euler步长。所有成功行都包含完整误差核算，以及附录~G
所述的另一数值构造检查。

\noindent\textbf{表~7. 四组Euler网格上的带符号 Heston 价格差证书。}
\begin{center}
\begin{tabularx}{\textwidth}{@{}l X X@{}}
\toprule
步长 & Euler价减连续价的包络 & 宽度上界 \\
\midrule
\(1/192\) & 无证书：投影条件检查失败 & 不可得 \\
\(1/384\) & \([-0.014169989568,0.013839530031]\) & 0.028009519599 \\
\(1/768\) & \([-0.011024692273,0.010642371599]\) & 0.021667063872 \\
\(1/1536\) & \([-0.010967952032,0.010766337995]\) & 0.021734290027 \\
\bottomrule
\end{tabularx}
\end{center}

在 \(h=1/192\) 时，固定的投影证明给出
\(\beta=184040303/23040000<8\)，而递推要求 \(\beta>8\)。
因此，这个充分条件不能认证该网格。三个成功行均满足绝对误差
容限 .025。表~8 分列其宽度贡献；未经舍入的精确条目之和
等于相应区间宽度。

\noindent\textbf{表~8. 三组成功网格中以价格单位计的宽度贡献。}
\begin{center}
\begin{tabularx}{\textwidth}{@{}X X X X@{}}
\toprule
宽度贡献 & \(1/384\) & \(1/768\) & \(1/1536\) \\
\midrule
收益函数余项 & 0.026754977094 & 0.020478587752 & 0.020171631177 \\
线性投影 & 0.000265787793 & 0.000008620825 & 2.65628E-7 \\
P 频率尾部 & 0.000028191423 & 0.000028191423 & 0.000028191423 \\
Q 频率尾部 & 0.000177646019 & 0.000368746602 & 0.000751284529 \\
周期化 & 0.000782917272 & 0.000782917272 & 0.000782917272 \\
舍入与端点 & 1E-12 & 1E-12 & 1E-12 \\
\bottomrule
\end{tabularx}
\end{center}

表~8 列的是宽度贡献，而非单侧误差半径；分别舍入后，显示值之和
可能有所变化。随网格细化，非线性部分占比分别约为 95.5\%、94.5\%
和 92.8\%。在最细步长处，离散频率尾部贡献的增加超过收益函数
余项和投影贡献的减少。(C.7) 中的频率尾界使用首步下界 \(i_0=hv_0\)，
因此，仅细化Euler网格不会收紧这一固定的反演方案。
这个宽度平台描述的是包络，而非真实价格误差。

% INSERT_STEPS_END

\subsection{第二个随机方差参数点}

保留合约、十二次月度观察、\(h=1/768\) 以及 (4.3) 中除初始方差外的
全部参数，将初始方差改为 \(v_0=.04\)。在该点重新计算加权矩、
投影前缀、变换，以及连续和离散尾部。同一增长估计给出
\(\exp(.6+4v_0)=\exp(.76)\)，因此保留的 \(\exp(.78)\) 矩界仍然保守。
两种数值构造均通过全部实载荷、分母、系数和复数分支条件检查。

\noindent\textbf{表~9. 初始方差为 .04 时的完整 Heston 证书。}
\begin{center}
\begin{tabularx}{\textwidth}{@{}l X@{}}
\toprule
量 & 经认证的包络或上界 \\
\midrule
连续模型价格 & \([6.499318797594,6.509948941256]\) \\
投影Euler价格 & \([6.498150434416,6.510702536853]\) \\
Euler价减连续价 & \([-0.011798506840,0.011383739259]\) \\
带符号区间宽度 & 0.023182246099 \\
\bottomrule
\end{tabularx}
\end{center}

该包络合并非线性余项、投影、两模型的频率尾部、周期化和向外算术。
六项精确宽度贡献与完整的另一构造检查，均记录在附录~G 的计算包中。
相较于 \(v_0=.045\)，带符号区间的宽度约增加
6.99\%。离散频率尾部贡献由约 .0003687466014 增至
.0016298806683，约占这一增量的 83.2\%。在 (C.7) 中，
较小的 \(v_0\) 降低首步下界 \(i_0=hv_0\)；首月拉普拉斯系数为负，
也使 \(\overline L_Q(U)=\exp(A+bv_0)\) 增大。这些因素使尾界变松。
该比较衡量包络的紧致程度，不表示真实偏差增大。
这是具有相同 Feller 指数的第二个点值证书。

\subsection{覆盖范围与敏感性设计}

主点的 Feller 指数为
\(2\kappa\bar v/\xi^2=2700/529\)，约为 5.104。
研究方差边界、相关性、执行价、初始方差和观察频率的敏感性，
须逐例提供新的矩、投影和尾部输入。特别是，改用 52 或 252 次观察
需要对齐的Euler网格，不能保留 \(h=1/768\)。首次观察时间也改变
收益函数界中的权重 \(I_1^{-1/2}\)。下述路径比较探索这些敏感性，
但不扩大经认证的参数范围。

\subsection{边界与观察频率敏感性的独立路径诊断}

路径实验每例使用 32,768 条路径，将一组Euler网格与步数加倍的网格
耦合。表~10 按 Feller 指数选取十四组实验中的三组；完整实验还改变
相关性、执行价、初始方差和步长。粗网格价减细网格价的估计具有
抽样误差，且保留细网格模型的离散化偏差。

\noindent\textbf{表~10. 三个 Feller 指数下的耦合路径诊断。}
\begin{center}
\begin{tabularx}{\textwidth}{@{}l X X@{}}
\toprule
情形 & 粗网格价减细网格价：估计值和标准误 & 粗网格投影比例 \\
\midrule
Feller 5.104 & -0.0000269; SE 0.0004717 & 0.0000000 \\
Feller 0.999 & 0.0001129; SE 0.0012453 & 0.0011268 \\
Feller 0.639 & 0.0009133; SE 0.0018375 & 0.0066991 \\
\bottomrule
\end{tabularx}
\end{center}

三个价格差估计的绝对值都小于一个估计标准误，因此表~10 既不能
判定其符号，也不能判定定价偏差是否系统性增加。提议值频率仅支持
如下观察：接近或低于边界时，负方差提议值更常出现；它们不是
投影误差界。

在 12、52 和 252 次观察下，\(I_1^{-1/2}\) 的样本均值分别为
16.489、34.078 和 74.871。对齐的粗网格分别有 768、832 和
1008 步，细网格步数加倍。这展示了平滑权重的敏感性，同时也
改变了时间网格。

\subsection{连续模型估值与数值风险预算}

假设估值系统为同一离散模型提供价格包络 \([q_L,q_U]\)。
定理~4.1 给出 \(e_h\in[\ell,u]\)，因此连续模型价格满足
\[
p_c\in[q_L-u,\ q_U-\ell].
\tag{8.1}
\]
在 (4.3) 下，直接得到的连续价格区间 (4.5) 宽度小于
\(.010498\)。这两种途径均为估值与实现检查提供以价格单位计的包络。

\subsection{校准后价格分布的稳定性}

数值实现的改变同时影响校准似然和目标价格。定理~7.1 保留两种影响：
\(\varepsilon_0+e_{P,J}+e_{Q,J}\) 控制固定参数下的目标价格差，
\(L_0\delta_u\) 控制参数后验分布移动引起的目标变化。

在原始 4096 单元网格上，输运贡献约占单报价界的 39\%，
占九报价界的 62\%。这促使我们进行下一小节的网格细化。
九报价包络更宽，反映的是较保守的输运贡献，并不意味着真实数值
偏差更大。

% INSERT_POSTERIOR_BEGIN
\subsection{后验网格细化}

表~11 保留相同的先验 (7.5)、合成观测、似然协方差和目标 Lipschitz
常数。每一行验证整个单元上的似然界与耦合不等式，并控制每个
亚式期权价格分位数的位移。

\noindent\textbf{表~11. 后验网格细化下亚式期权价格分位数位移的认证界。}
\begin{center}
\begin{tabularx}{\textwidth}{@{}l X X X@{}}
\toprule
单元数 & 单报价 & 九报价 & 耦合步数：单报价/九报价 \\
\midrule
4096 & 0.008018821658 & 0.012716404217 & 2 / 5 \\
8192 & 0.006452960804 & 0.010367612938 & 2 / 7 \\
16384 & 0.006061495591 & 0.009584682511 & 3 / 12 \\
\bottomrule
\end{tabularx}
\end{center}

共同的非网格误差贡献不变。细化后耦合步数增加，因此最终界不与
单元宽度严格同比缩放。16384 个单元下的新界比原始单报价和九报价界
分别约小 24.4\% 和 24.6\%。本次计算未评估自适应单元和局部
目标 Lipschitz 常数。

绝对分位数区间对应初始方差边缘分布的 2.5、50 和 97.5 百分位数。
对 4096 个单元下的连续模型后验，单报价区间约为
[.0307690,.0307764]、[.0444067,.0444141] 和 [.0590917,.0590992]；
九报价区间为 [.0308203,.0308350]、[.0435131,.0435425] 和
[.0587182,.0587256]。这些区间描述参数不确定性。尚未计算亚式期权
目标价格的绝对分位数，因此位移界还不能衡量这一变化相对于目标
后验不确定性尺度的大小。

% INSERT_POSTERIOR_END

\section{适用范围}

结构条件是认证的充分条件。第~5~节验证确定性方差类、指定的
Heston 模型及投影核；第~8~节改变主点步长，并认证第二个初始方差。
将这些包络扩展到参数邻域，需要一致的矩、变换、投影和尾部不等式。
第~7~节的小正 \(\xi\) 参数盒是另一个经认证的区域。

第~6~节给出 \(\xi=0\) 时的共同展开公式与显式余项界。
当 \(v_0\ne\bar v\) 时，九个看跌期权系数非零；亚式期权系数则由
得分积分和包含零的解析界给出。正 \(\xi\) 下的展开还需要额外
正则性，以及投影核和收益函数折点对应的可积余项。

后验位移界适用于受限先验 (7.5) 与合成报价。若要在五维参数域
\[
[2,4]\times[.03,.06]\times[.18,.28]\times[-.8,-.3]\times[.03,.06]
\]
上作迁移，其中坐标顺序为 \((\kappa,\bar v,\xi,\rho,v_0)\)，
需要一致的前向误差和后验输运输入。
亚式期权目标价格的绝对分位数，以及带符号的亚式期权首阶系数积分，
仍未计算。

分模型余项界导致第~8~节的宽度平台。命题 I.1 给出直接耦合余项不等式。
附录~I 的确定性方差例认证了所需的额外耦合矩，但尚未为随机方差
Heston 例建立这些矩。控制这些矩与收紧离散频率尾部，是使证书
宽度缩小所需的两个不同数学条件。

\section{结论}

共同高斯结构将经典对数正态条件化与投影Euler离散的随机方差模型
联系起来。核心引理把原始算术收益函数的折点余项转化为以积分方差
逆平方根加权的二阶矩。主定理将该量与有限变换证书合并，得到
带符号的价格差包络。指定 Heston 算例给出完整的货币单位误差界
\(.011025\)，以及经验证的连续模型价格区间。

步长网格实验表明，收益函数转换主导区间宽度，而固定的离散尾界
可能抵消细化Euler网格带来的收益。第二个随机方差点完成了针对
具体参数的完整执行；附录~I 的确定性例得到严格更窄的耦合非线性
区间。共同高斯展开与后验耦合还在各自所述参数族上提供公式层面
和价格单位下的界。
\appendix


\section{加权二阶矩的有限认证}

本附录以及附录 B、C 使用式 (4.3) 中的有理参数。记
\[
\alpha=\xi^2/2,\quad d=\kappa\bar v,\quad
s=1-\rho^2,\quad t_1=1/12.
\]
全文中，\(c\) 表示式 (4.3) 中的几何缩放系数。

\subsection{非负积分与载荷列表}

对 \(I>0\)，有
\[
I^{-1/2}=\frac2{\sqrt\pi}\int_0^\infty e^{-x^2I}\,dx.
\]
对非负平方项应用 Tonelli 定理，得到
\[
W_\nu=\frac2{\sqrt\pi}\int_0^\infty F_\nu(x)\,dx,\qquad
F_\nu(x)=E_\nu[(A-cG)^2e^{-x^2I}].
\tag{A.1}
\]
有 \(F_\nu\ge0\)，且该函数单调不增。在展开式 (5.2) 中，144 个有序的 \(A^2\) 项合并为 78 项；此外还有 12 个形如 \(S_iG\) 的项和一个 \(G^2\) 项。每个归一化指数项的股价总实载荷为 2，按月反向递推中的剩余载荷满足 \(p\in[0,2]\)。这 91 项的系数绝对值之和为
\[
S_0^2(1+c)^2=10000(1+c)^2.
\tag{A.2}
\]
首月施加杀灭率 \(x^2\)，其余月份的杀灭率均为零。

\subsection{真实连续流与离散代数递推}

在股价载荷为 \(q\)、杀灭率为 \(\lambda\) 的区间上，将归一化变换写成 \(\exp(\mathcal A+\mathcal Bv)\)。连续系数满足
\[
\mathcal B'=\alpha\mathcal B^2-(\kappa-\rho\xi q)\mathcal B+
(q^2-q)/2-\lambda,\qquad
\mathcal A'=rq+d\mathcal B.
\tag{A.3}
\]
在单步高斯积分中配方，得到离散代数递推
\[
\begin{aligned}
\mathcal A_j&=\mathcal A_f+h(rq+d\mathcal B_f),\\
\mathcal B_j&=\mathcal B_f+h\left[
\alpha\mathcal B_f^2-(\kappa-\rho\xi q)\mathcal B_f+
(q^2-q)/2-\lambda\right].
\end{aligned}
\tag{A.4}
\]
两式更新均使用更新前的 \(\mathcal B_f\)。附录 B 给出原始正部核与式 (A.4) 之间的差异界。固定观察日期通过按时间逆序加入对应的股价载荷来处理。

对实载荷 \(p\in[0,2]\) 及 \(\lambda\ge0\)，连续向量场在 \(\mathcal B=1\) 处严格向下。因此，从终值 0 出发的反向流满足 \(\mathcal B\le1\)，且 \(\mathcal A\le(2r+d)T\)。附录 B.3 中证明的一致可积性说明，该流就是实际的变换。

\subsection{经验证的 Darboux 包络}

取 \(\Delta x=1/2\)。连续模型使用节点 \(x_j=j/2,\ 0\le j\le256\)；离散模型使用 \(0\le j\le128\)。若真实节点值属于 \([f_j^-,f_j^+]\)，则
\[
\Delta x\sum_{j=1}^{M}\max(0,f_j^-)
\le\int_0^{M\Delta x}F_\nu(x)\,dx
\le\Delta x\sum_{j=0}^{M-1}\max(0,f_j^+).
\tag{A.5}
\]
这是非负、单调不增函数的矩形和界。对 \(Q\)，先将式 (A.2) 乘以附录 B 中每个指数项的误差界，再向两个方向扩大代数节点包络，最后应用式 (A.5)。

\subsection{离散模型的无穷尾部}

由于 \(I_Q\ge i_0=hv_0\)，对 \(x\ge X\)，有
\[
F_Q(x)\le F_Q(X)e^{-i_0(x^2-X^2)}.
\]
由 \(x^2-X^2\ge2X(x-X)\) 得到
\[
\int_X^\infty F_Q(x)\,dx\le\frac{\overline F_Q(X)}{2Xi_0},
\qquad X=64.
\tag{A.6}
\]
上界 \(\overline F_Q(X)\) 包含原始投影误差与舍入误差。

\subsection{连续模型的无穷尾部}

令 \(\kappa'=\kappa-2\rho\xi\)，并定义
\[
\Delta(x)=\sqrt{\kappa'^2+4\alpha(x^2-1)},\quad
r_1=\frac{\kappa'-\Delta}{2\alpha},\quad
r_2=\frac{\kappa'+\Delta}{2\alpha}.
\]
后续月份中的界 \(\mathcal B\le1\) 允许将首月流的终值取为 1，从而给出每个正载荷指数项的上界。对 \(x\ge X=128\)，有
\[
g=\frac{1-r_1}{1-r_2}\in(-1,0).
\]
Riccati 方程的交比解给出
\[
\mathcal B(t_1)\le r_1+\frac{\Delta}{\alpha}e^{-\Delta t_1},
\qquad
\int_0^{t_1}\mathcal B(t)\,dt
\le r_1t_1+\frac{\log2}{\alpha}.
\tag{A.7}
\]
此外，\(\Delta(x)\ge2\sqrt\alpha(x-1/X)\) 且 \(\Delta(X)t_1>1\)，故 \(\Delta e^{-\Delta t_1}\) 在整个尾部区域单调递减。定义
\[
C_0=\frac{\kappa'}{2\alpha}+\frac1{X\sqrt\alpha},\quad
C_1=C_0+\frac{\Delta(X)e^{-\Delta(X)t_1}}{\alpha},
\]
\[
C=\exp\!\left\{2r+d+2rt_1+
d\left(t_1C_0+\frac{\log2}{\alpha}\right)+v_0C_1\right\},
\qquad \zeta=\frac{dt_1+v_0}{\sqrt\alpha}.
\]
于是，每个归一化实变换均不超过 \(Ce^{-\zeta x}\)。由于 \((A-cG)^2\le2A^2+2c^2G^2\)，有
\[
\int_X^\infty F_P(x)\,dx
\le \frac{20000(1+c^2)C}{\zeta}e^{-\zeta X}.
\tag{A.8}
\]
所有根之间的关系、符号以及超越函数均采用向外舍入算术核查。

合并式 (A.5)--(A.8)，再乘以 \(2/\sqrt\pi\)，得到
\[
\begin{aligned}
W_P&\le2.172471592801589415440352,\\
W_Q&\le2.242189951599049647468543.
\end{aligned}
\tag{A.9}
\]
这些矩界同时证明了式 (2.7) 所要求的可积性。

\section{原始投影核与仿射表示的识别}

\subsection{单步投影修正}

设当前方差为 \(v\ge0\)，\(p\) 为股价载荷的实部。经过股价载荷诱导的高斯倾斜后，下一步候选方差满足
\[
Y\sim N(\eta_pv+dh,\;\xi^2hv),\qquad
\eta_p=1-(\kappa-\rho\xi p)h,
\]
外部因子为 \(\exp(rph+\gamma_p hv)\)，其中 \(\gamma_p=(p^2-p)/2\)。对后续复方差系数 \(b\)，若候选值为 \(Y=-z<0\)，则
\[
|e^{bY^+}-e^{bY}|
\le |b|z e^{|b|z}
\le \frac{|b|h}{\lambda e}
\exp[(\lambda/h+|b|)z].
\tag{B.1}
\]
令 \(s_b=\lambda+h|b|\)。由高斯矩公式，依赖状态的误差界中的指数化为
\[
-\frac v h\left(s_b\eta_p-\alpha s_b^2-\gamma_ph^2\right)-s_bd.
\tag{B.2}
\]
取 \(\lambda=12\)。假设 \(p\) 位于指定载荷范围，且 \(\eta_p\ge\eta_*\)、\(\gamma_p\le\gamma_*\)、\(|b|\le b_*\)。只需检查该凹二次式在两个端点 \(s_b=12\) 和 \(s_b=12+b_*h\) 处的下界，就能保证括号中的表达式至少为 8。

\subsection{从局部误差界到真实前段路径}

定义
\[
u_0=8,\qquad u_{j+1}=\eta_*u_j-\alpha u_j^2-\gamma_*h^2.
\tag{B.3}
\]
在全部 768 个层级检查非负性 \(u_j\ge0\)。对 \(u\ge0\)，有 \(e^{-uY^+/h}\le e^{-uY/h}\)。因此，反复应用高斯矩界即可给出原始核下股价加权前段路径的拉普拉斯期望上界。由此得到每个完整指数项的投影误差界：
\[
\delta_Q=
\frac{b_*}{12e}\exp(C_f+C_{\rm pre}-12d)\,
h\sum_{j=0}^{767}
\exp\!\left[-d\sum_{k<j}u_k-\frac{v_0u_j}{h}\right].
\tag{B.4}
\]
这里，\(C_f\) 给出后续常数项的上界，\(C_{\rm pre}\) 给出真实前段路径上股价漂移的上界。两组参数如下。

\noindent\textbf{表~B.1. 投影误差预算的常数。}

\begingroup
\small
\setlength{\tabcolsep}{3pt}
\renewcommand{\arraystretch}{1.2}
\begin{tabularx}{\textwidth}{@{}>{\raggedright\arraybackslash}X c r c c c@{}}
\toprule
用途 & \(p\) 的范围 & \(b_*\) & \(\eta_*\) & \(\gamma_*\) & \(C_f+C_{\rm pre}\) \\
\midrule
加权二阶矩 & \([0,2]\) & 512 & \(1-(\kappa-2\rho\xi)h\) & 1 & \(4r+d\) \\
十三个线性变换 & \([-1/2,1]\) & 1024 & \(1-(\kappa-\rho\xi)h\) & \(3/8\) & \(2r+d\) \\
\bottomrule
\end{tabularx}
\endgroup

有限列表同时检查 \(\operatorname{Re}b\le1\) 及后续系数的模长界。第一种情况下，将式 (B.4) 乘以 \(10000(1+c)^2\)，得到每个拉普拉斯节点的误差预算。第二种情况下，将它乘以所有傅里叶系数的绝对值之和；此和不再额外乘以 \(h\)。

\subsection{矩增长与停止操作的去除}

对实股价载荷 \(p\)，取 \(e^{4V}\) 作为测试函数的方差部分。连续生成元计算中 \(V\) 的系数为
\[
16\alpha-4\kappa+4\rho\xi p+\gamma_p.
\tag{B.5}
\]
对所有所需载荷，该系数均严格为负，包括按 \(5/4\) 倍缩放后的范围 \(p\in[-5/4,5/2]\)。对 \(Q\)，正系数 4 允许直接使用 \(1-e^{-4z}\le4z\)。结合 \(\lambda=12\) 时的高斯尾界，投影引起的额外增长率至多为
\[
\frac{4e^{-12d}}{12e}.
\]
在同一扩大的范围内，倾斜指数中的状态系数仍非正。例如，使用最不利值 \(\eta_{5/2}\) 和 \(\gamma_*=15/8\)，精确计算得到
\[
12\eta_{5/2}-144\alpha-\frac{15}{8}h^2
=\frac{8001336523}{983040000}>0.
\]
总常数增长满足
\[
4d+\frac52r+\frac{4e^{-12d}}{12e}
<.589267620943<.6.
\tag{B.6}
\]
因此，真实连续模型与离散模型所需的指数矩均有一致的有限上界。特别地，对原始载荷范围 \(p\in[-1,2]\)，归一化矩的上界为
\[
M_2=e^{4v_0+.6}=e^{.78}
\tag{B.7}
\]

先将连续仿射表达式停止在有界方差区域内。将其模长取 \(5/4\) 次幂，会把股价载荷缩放到上述范围。后续正方差系数至多为 \(5/4<4\)，而非正的杀灭项可在求上界时略去。因此，式 (B.5)--(B.6) 对停止后的族给出一致的 \(L^{5/4}\) 界，从而得到一致可积性。去除停止操作后，Riccati 表达式即为真实变换的期望。有限离散望远镜和只包含有限个可积的条件期望，同一矩增长估计适用于这些期望。

\section{十三个线性变换反演的完整误差控制}

Lee~\cite{Lee2004} 分析了傅里叶定价中的截断误差与采样误差。下面的界使用本文十三个指数项的高斯衰减及周期化矩。

\subsection{阻尼反演与有限和}

对正测度 \(\mu\)，记 \(\widehat\mu(z)=\int e^{zy}\mu(dy)\)。取
\[
\begin{gathered}
\delta=1/2,\qquad \Delta u=1/5,\qquad N=640,\\
U=N\Delta u=128,\qquad L=2\pi/\Delta u=10\pi.
\end{gathered}
\]
当负阶矩有限时，\(e^{-\delta y}F_\mu(y)\) 可积，其傅里叶变换为 \(\widehat\mu(-\delta-iu)/(\delta+iu)\)。因此，
\[
F_\mu(y)=\frac1{2\pi}\int_{\mathbb R}
\frac{e^{(\delta+iu)y}}{\delta+iu}\widehat\mu(-\delta-iu)\,du.
\tag{C.1}
\]
代入式 (3.6)，令 \(z_n=\delta+in\Delta u\)、\(w_0=1/2\)，以及 \(w_n=1\)（当 \(n\ge1\) 时）。有限实和中的系数为
\[
a_{0,n}=\frac{e^{-r}\Delta u\,w_n}{\pi z_n}
\left(95e^{z_ny_{95}}-110e^{z_ny_{110}}\right),
\]
\[
a_{i,n}=\frac{e^{-r}\Delta u\,w_n}{\pi z_n}\frac{100}{12}
\left(e^{z_ny_{110}}-e^{z_ny_{95}}\right),\qquad 1\le i\le12.
\tag{C.2}
\]
对应的变换为 \(Ee^{-z_nY}\) 与 \(Ee^{Z_i-z_nY}\)。零频率取半权重，而保留的最高频率取完整权重。股价指数项中剩余载荷的实部位于 \([-1/2,1]\)；其模长平方需要范围 \([-1,2]\) 内的实载荷。

\subsection{周期化误差}

Poisson 周期化给出
\[
\sum_{k\in\mathbb Z}e^{-\delta kL}F_\mu(y+kL).
\]
对 \(k>0\)，使用总质量 \(M_{\mu,0}\)；对 \(k=-j<0\)，使用 \(F_\mu(y-jL)\le e^{y-jL}M_{\mu,-1}\)。于是，非零指标各项之和至多为
\[
\frac{M_{\mu,0}}{e^{\delta L}-1}
+\frac{M_{\mu,-1}e^y}{e^{(1-\delta)L}-1}.
\tag{C.3}
\]
这里，\(M_{0,0}=1\)、\(M_{A,0}\le100e^{.01}\)、\(M_{0,-1}\le e^{.78}\)、\(M_{A,-1}\le100e^{.78}\)。在价格组合中逐项取绝对值，得到单个概率律的周期化误差预算：
\[
E_{\rm alias}=
e^{-.01}\frac{
205+200e^{.01}
+e^{.78}(195e^{y_{95}}+210e^{y_{110}})
}{e^{5\pi}-1}.
\tag{C.4}
\]
两个模型的差异计入 \(2E_{\rm alias}\)。下述频率尾界保证傅里叶系数绝对可和，而式 (C.3) 保证周期化表达式在紧区间上绝对且一致收敛。因此，傅里叶表示的唯一性说明，两种表示对应同一连续函数。

\subsection{共同高斯因子引起的频率衰减}

在式 (5.3) 的条件化下，每个指数项作用于首月因子 \(U\) 的总虚载荷为 \(-u\)。令 \(K=E[e^{\text{real loading}}\mid\mathcal G]\)。先作条件高斯积分，再使用 Cauchy--Schwarz 不等式和条件 Jensen 不等式，得到
\[
|\Phi_\nu(u)|
\le E[K e^{-u^2sI/2}]
\le\sqrt{E e^{2\,\text{real loading}}}
\sqrt{E e^{-su^2I}}
\le\sqrt{M_2L_\nu(u)}.
\tag{C.5}
\]
这里，\(L_\nu(u)=E_\nu e^{-su^2I}\)。对定义 \(F_A\) 的十二项平均值，保留额外的价格水平因子 100。定义
\[
C_*=
e^{-.01}\left(195e^{\delta y_{95}}+210e^{\delta y_{110}}\right).
\]

在 \(Q\) 下，\(I\ge i_0=hv_0\)，故
\[
L_Q(u)\le L_Q(U)e^{-si_0(u^2-U^2)}\quad(u\ge U).
\]
真实 \(L_Q(U)\) 的经验证上界仅需 64 步反向递推。从 \(b=A=0\) 出发，使用
\[
A\leftarrow A+dhb,\qquad
b\leftarrow b+h(\alpha b^2-\kappa b-sU^2).
\tag{C.6}
\]
每一层均检查 \(b\le0\)。不等式 \(e^{bY^+}\le e^{bY}\) 随后保证，该高斯递推给出原始正部核的上界。最终，\(\overline L_Q(U)=e^{A+bv_0}\)。用积分控制正递减包络的离散和，得到
\[
E_{{\rm tail},Q}\le
\frac{C_*\sqrt{M_2\overline L_Q(U)}}{\pi s i_0U^2}.
\tag{C.7}
\]
这里使用的积分不等式为
\[
\int_U^\infty e^{-a(u^2-U^2)}\frac{du}{u}\le\frac1{2aU^2}.
\]

对 \(P\)，定义
\[
\Delta_U=\sqrt{\kappa^2+4\alpha sU^2},\quad
C_0=\kappa/(2\alpha),\quad
C_1=C_0+\Delta_Ue^{-\Delta_Ut_1}/\alpha,
\]
\[
C_P=\exp\!\left\{
d(t_1C_0+\log2/\alpha)+v_0C_1\right\},\qquad
\gamma=\frac12\sqrt{s/\alpha}(dt_1+v_0).
\]
采用式 (A.7) 中相同的 CIR 交比界，但令终值为 \(b=0\)，得到 \(L_P(u)\le C_Pe^{-2\gamma u}\)，适用于每个 \(u\ge U\)。因此，
\[
E_{{\rm tail},P}\le
\frac{C_*\sqrt{M_2C_P}}{\pi\gamma U}e^{-\gamma U}.
\tag{C.8}
\]

\subsection{连续复流的分支}

记 \(\kappa_q=\kappa-\rho\xi q\)、\(\gamma_q=(q^2-q)/2\)、\(D_q=\sqrt{\kappa_q^2-4\alpha\gamma_q}\)。将每个月划分为八个长度为 \(\Delta t=1/96\) 的子区间。M\"obius 流的分母为
\[
\mathcal D_t=\cosh(D_qt/2)
+(\kappa_q-2\alpha b)\frac{\sinh(D_qt/2)}{D_q}.
\]
该流的积分为
\[
\int_0^{\Delta t}b(t)\,dt
=\frac{\kappa_q\Delta t/2-\log\mathcal D_{\Delta t}}{\alpha}.
\tag{C.9}
\]
在每个子区间上检查
\[
\cosh(|D_q|\Delta t/2)-1+
\frac{|\kappa_q-2\alpha b|\Delta t}{2}
\cosh(|D_q|\Delta t/2)<1.
\tag{C.10}
\]
幂级数的绝对值界说明，在整个子区间上有 \(|\mathcal D_t-1|<1\)。因此，主值对数始终与分母的初值 1 连续相连。在边界 \(\operatorname{Re}b=1\) 上，虚部二次型
\[
-\alpha(\operatorname{Im}b)^2
-\rho\xi(\operatorname{Im}q)(\operatorname{Im}b)
-(\operatorname{Im}q)^2/2
\]
非正，且其余实漂移严格为负。因此，连续流满足 \(\operatorname{Re}b\le1\)。附录 B.3 说明，该仿射表达式就是真实的变换。

\subsection{误差合并与精度}

令 \(\widehat e_L\) 为有限代理差。真实线性差的包络为
\[
\widehat e_L+
[-E_L,E_L],\qquad
E_L=2E_{\rm alias}+E_{{\rm tail},P}
+E_{{\rm tail},Q}+\delta_Q\sum_{n,i}|a_{i,n}|.
\tag{C.11}
\]
所有有限和与系数均计算为 Arb 包络。代入精确有理端点得到
\[
E_L\le.000594238059839992574505.
\]
将此结果与式 (A.9) 给出的非对称收益余项合并，并按定理 4.1 处理，即得式 (4.4)。


\section{高斯密度展开与有限步长界}

\subsection{方差展开}

对 \(t/h\in\mathbb N\)，令
\[
x=-\frac th\log(1-\kappa h)-\kappa t\ge0.
\]
对数的幂级数给出
\[
0\le x-\frac{\kappa^2th}{2}
\le\frac{\kappa^3th^2}{3(1-\kappa h)},\qquad
x\le\frac{\kappa^2th}{2(1-\kappa h)}.
\]
使用 \(|1-e^{-x}-x|\le x^2/2\)，并代入
\[
I_h(t)-I(t)=\frac D\kappa e^{-\kappa t}(1-e^{-x})
\]
得到 \(|I_h(t)-I(t)-ha(t)|\le h^2B_I(t)\)。
将相邻时点的两个累计余项界相加，即得式 (6.5)。

\subsection{密度导数}

对单个增量，令 \(X=Y-r/12\)。其密度为
\[
f_s(y)=(2\pi s)^{-1/2}
\exp[-X^2/(2s)-X/2-s/8].
\]
因此，
\[
\partial_s\log f_s=-\frac1{2s}+\frac{X^2}{2s^2}-\frac18,\qquad
\partial_s^2\log f_s=\frac1{2s^2}-\frac{X^2}{s^3}.
\]
在 \(X=-s/2+\sqrt{s}G\) 下，一阶得分函数的均值为零，二阶矩为
\(J(s)=1/(2s^2)+1/(4s)\)。此外，
\[
E|\partial_s^2\log f_s|
\le\frac3{2s^2}+\frac1{4s}.
\]
各坐标的独立性使交叉得分的期望为零，由此得到两个 \(L^1\) 导数界，用于定理~6.1 的证明。
在具有正下界和有限上界的方差盒上，一阶与二阶密度导数均一致地受
\(C(1+|y|^4)e^{-c|y|^2}\) 控制。
因此，可以按需交换差商、参数积分与收益积分。

\subsection{完整的有限步长比较}

对两个方差向量 \(s,\tilde s\)，假设端点及连接它们的线段均满足 \(s_j,\tilde s_j\ge u_j>0\)。
沿线段积分一阶得分函数，得到
\[
\operatorname{TV}(p_s,p_{\tilde s})
\le\left\{
\sum_j(\tilde s_j-s_j)^2
\left(\frac1{8u_j^2}+\frac1{16u_j}\right)
\right\}^{1/2}.
\tag{D.1}
\]
对值域长度为 \(L\) 的收益，价格差至多为
\(Le^{-r\tau}\operatorname{TV}\)。
在式 (6.10) 的范围内，使用一致界 \(|D|\le.015\) 和 \(u_j\ge.03/12\)，有限计算得到
\[
|e_{h,\rm Asian}|\le.004990923469,\qquad
|R_{h,\rm Asian}|\le.000165216416.
\]
后验模块使用每个月更紧的实际最小方差，因此其在固定参数下的参考界 \(\varepsilon_0\) 可以小于上述一致界。

\subsection{首阶系数的有效可计算表示}

对有理参数和本文所考虑的显式 Lipschitz 收益，式 (6.7) 中的首阶系数是有效可计算的积分。令
\[
S_d(G)=\sum_jd_j\left[\frac{G_j^2-1}{2s_j}-\frac{G_j}{2\sqrt{s_j}}\right],
\qquad F=\sum_jd_j^2J(s_j).
\]
将收益中心化后，略去区域 \(\max_j|G_j|>R\) 引起的误差至多为
\[
e^{-r\tau}\frac L2\sqrt{24F}\,e^{-R^2/4}.
\tag{D.2}
\]
这直接来自 Cauchy--Schwarz 不等式和高斯尾概率的并集界。
在 \([-R,R]^{12}\) 上，定义
\[
\begin{aligned}
M_S&=\sum_j|d_j|\left[\frac{R^2+1}{2s_j}+\frac R{2\sqrt{s_j}}\right],\\
L_S&=\sum_j|d_j|\left[\frac R{s_j}+\frac1{2\sqrt{s_j}}\right].
\end{aligned}
\]
若收益作为观察时点对数股价向量的函数，其 Lipschitz 常数为 \(L_z\)，则被积函数的 Lipschitz 常数至多为
\[
L_H=L_z\left(\sum_j\sqrt{s_j}\right)M_S+\frac L2L_S.
\]
对边长为 \(\ell\) 的各个盒子，将其中点值乘以精确高斯盒概率后求和，总内部误差至多为
\(e^{-r\tau}L_H\ell/2\)。
选取 \(R,\ell\)，并结合式 (D.2)，便得到最多需要
\(\lceil2R/\ell\rceil^{12}\) 个盒子的有限算法。
这一表示建立了可计算性；第~9 节说明，它并不等同于已完成首阶系数的高精度积分计算。

\section{正波动率之波动率下的扰动常数}

在连续模型中，\(D_t=V_t^\xi-v(t)\) 具有线性常数变易表示。由于 \(EV_t^\xi=v(t)\le M\)，It\^o 等距关系给出
\[
ED_t^2=\xi^2\int_0^te^{-2\kappa(t-s)}EV_s^\xi\,ds
\le\frac{\xi^2M}{2\kappa_-}.
\tag{E.1}
\]
离散确定性参考方差保持非负。
正部映射是 1-Lipschitz 的，高斯创新的条件均值为零。因此，记 \(D_j^{(2)}=E(V_j^\xi-v_j)^2\)，得到
\[
D_{j+1}^{(2)}
\le(1-\kappa h)^2D_j^{(2)}+\xi^2hEV_j^\xi
\le(1-\kappa h)^2D_j^{(2)}
+\xi^2h(M+\sqrt{D_j^{(2)}}).
\tag{E.2}
\]
取 \(C_Q^2=2M/\kappa_-\)。
条件 (7.1) 给出 \(\xi C_Q\le M\) 及
\(1-(1-\kappa h)^2\ge\kappa_-h\)。
从零初始误差出发归纳，得到
\(D_j^{(2)}\le\xi^2C_Q^2\)。
这一步保留了投影对 \(EV_j^\xi\) 的影响。

对确定性参考值 \(v\ge m\)，有
\[
|\sqrt{x}-\sqrt v|^2
=\frac{|x-v|^2}{(\sqrt x+\sqrt v)^2}\le |x-v|^2/m.
\]
对数股价差中的漂移贡献以 \(T\xi C_\ell/2\) 为界。
It\^o 等距关系和 Doob 的 \(L^2\) 不等式将鞅贡献控制在
\(2\xi C_\ell\sqrt{T/m}\) 以内，从而证明式 (7.3)。
股价布朗运动与方差驱动之间的相关性，不改变这些针对适应随机积分的不等式。

对看跌期权，在收益非零区域内，收益对对数股价的导数绝对值至多为 \(K\)。
对封顶亚式收益，考虑观察时点对数股价向量中的一条线段。
在收益变化区域 \(K_1<A<K_2\) 中，梯度的 \(\ell^1\) 范数为 \(A\le K_2\)；其余区域为零。
分段绝对连续性给出全局 Lipschitz 常数 110。
因此，对到期时间 \(T\)，有
\[
e_{\ell,i}=K_i e^{-rT_i}\xi_{\max}C_\ell
\left(T_i/2+2\sqrt{T_i/m}\right),
\]
\[
e_{\ell,J}=110e^{-r}\xi_{\max}C_\ell
\left(1/2+2/\sqrt m\right).
\tag{E.3}
\]
在式 (7.5) 中，\(m=.03,M=.06,\kappa=3\)，因此 \(C_P=.1,C_Q=.2\)，得到
\[
e_{P,J}\le.000131198492451391,\qquad
e_{Q,J}\le.000262396984902781.
\tag{E.4}
\]
对式 (7.4)，先采用向外舍入算术，在 \(\kappa\) 的 256 个闭区间上包络参考方差差异，再应用
\[
|p_h^\xi-p_c^\xi|
\le |p_h^0-p_c^0|+
110\xi(C_P+C_Q)(1/2+2/\sqrt m)
\]
得到正文给出的整个参数盒价格界。
其余参数通过对
\(m,M,|v_0-\bar v|\) 的一致界予以包络，而不是用端点样本代替完整包络。

\section{后验输入、归一化与有限耦合}

\subsection{精确报价}

报价按 \(T=1/4,1/2,1\) 排列，每个到期时间内按 \(K=90,100,110\) 排列。
下列有理数作为给定的合成输入：
\[
y=\begin{pmatrix}
526706580970763/562949953421312\\
285901443526503/70368744177664\\
754049251176359/70368744177664\\
580752698287259/281474976710656\\
794307614955649/140737488355328\\
830842010134081/70368744177664\\
269207942741719/70368744177664\\
137257851354697/17592186044416\\
1918003280297167/140737488355328
\end{pmatrix}.
\tag{F.1}
\]
这些分数是参考输入文件中 binary64 数值的精确值；单报价情形使用第八个条目。
它们是似然函数的输入，生成这些数值的辅助计算不应被视为连续模型价格的证明。

\subsection{参考价格与目标常数}

令
\[
b_P(t)=\frac{1-e^{-3t}}3,\quad
b_Q(t)=\frac{1-(1-3h)^{t/h}}3,
\]
\[
I_m(T;u)=.045T+(u-.045)b_m(T).
\]
参考看跌期权价格通过精确的 Black 公式计算：
\[
\begin{aligned}
p_m(K,T;u)&=Ke^{-rT}\Phi(-d_2)-S_0\Phi(-d_1),\\
d_1&=\frac{\log(S_0/K)+rT+I_m/2}{\sqrt{I_m}},\qquad
d_2=d_1-\sqrt{I_m}.
\end{aligned}
\tag{F.2}
\]
由于 \(b_m(T)>0\) 且 vega 为正，整个 \(u\) 单元上的参考价格可由其端点值包络。
再按式 (E.3) 扩大两端点，就能包络整个单元及所有 \(\zeta\) 上的真实价格。

令
\[
\begin{aligned}
a_{m,j}&=b_m(t_j)-b_m(t_{j-1}),\\
z_{m,j}(u)&=.045/12+(u-.045)a_{m,j},\qquad B_*=15e^{-.01}.
\end{aligned}
\]
附录~D 中的中心化得分界给出
\[
L_0=\frac{B_*}2
\sqrt{\sum_ja_{P,j}^2J(z_{P,j}(.03))},
\]
\[
\begin{aligned}
\varepsilon_0&=\frac{B_*}2
\sqrt{\sum_j\Delta_j^2J(\underline z_j)},\\
\Delta_j&=.015|a_{Q,j}-a_{P,j}|,\qquad
\underline z_j=\min(z_{P,j}(.03),z_{Q,j}(.03)).
\end{aligned}
\tag{F.3}
\]

\subsection{整个单元上的似然包络}

对 \(d=1\) 或 9，有
\[
\Sigma_d^{-1}=\frac{64}{27}\left(I_d-\frac{\mathbf1\mathbf1^T}{d+3}\right).
\]
将残差向量记为 \(r\)，其势函数是非负项之和：
\[
\frac12r^T\Sigma_d^{-1}r=
\frac{32}{27(d+3)}
\left(3\sum_i r_i^2+\sum_{i<j}(r_i-r_j)^2\right).
\tag{F.4}
\]
对该表达式及其指数函数采用向外舍入计算，再舍入到间距为 \(2^{-96}\) 的精确整数格上，得到
\[
0<l_{m,i}\le w_m(u,\zeta)\le u_{m,i}\le1
\]
该界在整个单元与干扰参数域上成立。
对概率测度 \(\nu(d\zeta)\) 积分，可得到 \(\bar w_m(u)\) 的同一包络。

\subsection{共享归一化与逆 CDF 配对}

每个单元具有相同的先验质量。
在节点 \(x_k\) 处，将左侧质量下界与上界之和记为
\(L_<,U_<\)，右侧相应的和记为 \(L_>,U_>\)。
边际 CDF \(G_m\) 满足
\[
\underline G_m(x_k)=\frac{L_<}{L_<+U_>},
\qquad
\overline G_m(x_k)=\frac{U_<}{U_<+L_>}.
\tag{F.5}
\]
函数 \(A/(A+B)\) 关于 \(A\) 递增、关于 \(B\) 递减，因此这些包络保留了分子和分母中的共享质量。
归一化常数也直接通过对所有权重求和得到包络。

令 \(U\) 为共同的均匀随机变量，并置 \(U_m=G_m^{-1}(U)\)。
若 \(U_m\in[x_i,x_{i+1}]\)，则
\[
U\in[\underline G_m(x_i),\overline G_m(x_{i+1})].
\tag{F.6}
\]
对每个 P 单元，枚举所有 Q 单元区间中最左和最右的指标 \(j_-,j_+\)；这些区间在式 (F.6) 中对应的范围可能与该 P 单元的范围相交。
令
\[
\delta_u=\Delta x\max_i\max\{|i-(j_++1)|,\ |i+1-j_-|\},
\quad \Delta x=.03/4096.
\tag{F.7}
\]
真实单元对必然属于该候选集合，因而
\(|U_P-U_Q|\le\delta_u\) 几乎处处成立。
所有比较均使用精确整数和分数。

从所有单元得到的真实三维先验归一化常数包络见表~F.1。

\noindent\textbf{表~F.1. 4096 个单元上的后验归一化常数包络。}
\begin{center}
\begin{tabularx}{\textwidth}{@{}l X X@{}}
\toprule
情形 & \(Z_P\) & \(Z_Q\) \\
\midrule
单报价 &
\(\begin{gathered}[c] [.935743002637584884,\\ .935913179433726711]\end{gathered}\) &
\(\begin{gathered}[c] [.935667223967444023,\\ .935927287640792653]\end{gathered}\) \\
\addlinespace
九报价 &
\(\begin{gathered}[c] [.620128921272339539,\\ .621155228251382142]\end{gathered}\) &
\(\begin{gathered}[c] [.619743550733119263,\\ .621256330676942526]\end{gathered}\) \\
\bottomrule
\end{tabularx}
\end{center}

将附录~E 的目标半径与式 (F.3)、(F.7) 合并，即得式 (7.8)。
该过程包络每个模型各自的完整后验。

\section{可复现实现}

从以下地址的 \texttt{paper} 标签开始：
\url{https://github.com/130U/certified-valuation-arithmetic-asian-options/tree/paper}。
其中包含稿件、步长网格和后验网格实验、第二个 Heston 参数点以及确定性耦合例子。
根目录中的 \texttt{REVISION-MANIFEST.json} 列出了该检出版本的文件及其哈希；
\texttt{code/MANIFEST.json} 保留了冻结原始数值包的身份。

历史数据仍可在以下版本中获取：网格实验见
\href{https://github.com/130U/certified-valuation-arithmetic-asian-options/tree/cf0d242f1590c7ca07ab8ca7afb57da5a583032b}{提交 \texttt{cf0d242f1590c7ca07ab8ca7afb57da5a583032b}}，
第二个参数点及耦合数据见
\href{https://github.com/130U/certified-valuation-arithmetic-asian-options/tree/1959f8f078e1fedd590000d6166cb20429ff31db}{提交 \texttt{1959f8f078e1fedd590000d6166cb20429ff31db}}。
原始核的基线版本为 \texttt{5bc72cf6036fdd73fea9c8bc6c85f85fd474b79a}。

\noindent\textbf{表~G.1. 数值模块及其结果。}
\begin{center}
\begin{tabularx}{\textwidth}{@{}l X@{}}
\toprule
模块 & 对应结果 \\
\midrule
\texttt{asian} & 原始参数点上的定理 4.1 与推论 4.2。 \\
\texttt{tt} & 定理 6.1 与确定性参数面 (6.10) 上的常数。 \\
\texttt{small-xi} & 五维扰动参数盒 (7.4)。 \\
\texttt{posterior-zero} & 高斯参考后验及目标常数。 \\
\texttt{posterior} & 受限先验 (7.5) 与定理 7.1。 \\
\bottomrule
\end{tabularx}
\end{center}

参考运行环境是未启用优化的 CPython 3.12，运行于 Windows x86-64，并配合 \texttt{python-flint==0.8.0}。
在发表版本的检出目录中，下列命令创建环境、核验原始包，并使用两套实现执行原始亚式证书。
每个双连字符均为命令行中的字面字符。

\begin{verbatim}
py -3.12 -m venv .venv
./.venv/Scripts/python.exe -m pip install -r code/requirements.txt
./.venv/Scripts/python.exe code/run.py verify
./.venv/Scripts/python.exe code/run.py run --module asian --independent
\end{verbatim}

运行会打印其在 \texttt{code/runs/} 下的新目录。
将下面的 \texttt{run-ID} 替换为该目录名。
若要执行所有原始模块，将上一条命令中的 \texttt{asian} 替换为 \texttt{all}。

\begin{verbatim}
./.venv/Scripts/python.exe code/run.py check --run code/runs/run-ID
\end{verbatim}

安装可能需要联网；数值计算无需联网。
包内不含依赖 wheel 文件，离线安装需要另行取得可信且兼容的 wheel。
断言约束数学条件与资源契约。
每个工作进程使用一个线程，硬内存上限为 256 MiB。

原始完整执行及新增实验保存在 \texttt{code/revision/results/evidence.zip} 中。
其清单绑定每个成员的字节与 SHA-256。
索引 \texttt{code/revision/results/INDEX.md} 提供档案提取与账本复核命令；
\texttt{code/revision/README.md} 提供重新计算命令及实验输入。
已保存结果的核验与重新执行数值核是两种不同操作。

步长网格账本 \texttt{code/revision/results/revision-summary.json} 将每个印出的数值行对应到精确源字段、结果、哈希与收据。
原始加权矩保存在 \texttt{code/reference/asian-remainder-result.json}；
线性中心、变换贡献及有符号价格包络保存在 \texttt{code/reference/asian-linear-result.json}。
精确合成观测保存在 \texttt{code/data/synthetic\_quotes.json}。
附录 A、B、C 分别说明非线性项列表、原始投影递推与线性反演。

参考比较包括所有数值 JSON 字段，包括数组和有理端点。
仅排除哈希元数据及六个具名耗时字段。
下述替代构造共享 Arb 和解析不等式；数值一致性验证的是这些假设下已实现的有限计算，并不提供对共同不等式的独立推导。

实验档案同时保留中止尝试与成功完成的运行。
在 192 步时，固定投影不等式未满足所要求的指数界。
最初的 384 步尝试在原始数值半径目标处停止；成功的包络运行将该目标记录为 false，同时保留所有解析条件。
在 1536 步时，最初工作进程超过了原有墙钟时间预算。
完成的运行将该有限预算乘以 \texttt{ceil(N/768)}，保留数学检查、精度、线程数及内存限制。
另行开展的加权矩实验使用同样类型的已声明运行时间缩放。
收据区分解析条件排除、数值验收目标与资源中止。
所报告的任务耗时取决于机器及负载。


\subsection{数值构造与已记录成本}


在主参数点，非线性矩模块计算 24,388 个连续实流块和 815,360 个离散实递推步。
线性模块计算 641 个频率和十三个指数项，共执行 6,399,744 个离散步及 799,968 次连续分支检查。
主程序使用 256 位 Arb~\cite{Johansson2017}；核查程序采用 384 位 Arb，并使用不同的数值构造。

对加权余项，\texttt{check-asian-remainder.py} 将主程序的根与交比 Riccati 公式替换为双曲二乘二流，通过有序数对构建载荷列表，并单独索引反向离散网格。
它将投影前段路径因子累积为乘积，而不是对一个和取指数。
对线性反演，\texttt{check-asian-linear.py} 将主程序的实部、虚部分开递推替换为复数 \texttt{acb} 递推，并用指数函数表示连续双曲块。
它核查每个频率贡献，并在不调用主程序数值函数的情况下重新计算有限误差预算。
两套核查共享 Arb、解析恒等式与理论界，检验的是实现一致性和算术稳定性，并不提供这些界的独立推导。
根流和矩阵流计算也包络了式 (4.11) 使用的同一组分组正矩。

原始非线性模块和线性模块的历史工作进程耗时分别约为 12.02 秒和 33.20 秒，每个进程使用一个线程及 256~MiB 硬内存上限。
在步长网格研究中，五个亚式任务的外部墙钟总时间分别约为 122.3、184.0、431.6 秒，对应 384、768、1536 步。
这些后续观测包括记录负载下的控制器及核查程序，不能作为可移植的性能保证。
每份收据保留任务耗时、操作计数与资源契约。



\subsection{抽样示例与后验边际包络}


亚式系数示例使用 binary64 算术，每个参数点使用 2,097,152 个高斯增量向量。
它计算精确确定性方差高斯系数被积函数和密度比余项，四个步长共享同一组向量。
减去已知均值为零的得分项，并使用 \texttt{expm1}，可减少缩放后余项中的抵消损失。
输入、种子、公式、抽样估计与估计标准误保存在 \texttt{code/revision/results/asian-beta-diagnostic.json} 中。
其正态近似抽样区间不同于第~6 节中的解析系数包络及确定性余项界。

另行开展的 NumPy 路径研究对十四种情形中的每一种使用 32,768 条路径，将原始投影 Euler 网格与步数加倍的网格耦合。
文件 \texttt{code/revision/results/path-diagnostics.json} 记录种子、binary64 输入、抽样价格、标准误和负候选值计数。
细网格概率律仍有离散偏差，因此这些抽样结果不包络连续模型价格。

后验网格汇总 \texttt{code/revision/results/posterior-grid-summary.json} 保留完整有理 CDF 包络，以及初始方差边际分布的 2.5、50、97.5 百分位分位数括区间。
绝对亚式价格目标分位数不属于该计算的内容。


\subsection{第二个参数点与耦合矩的复算}

第~8.3 节完成的第二个 Heston 参数点，以及附录 I 中已执行的确定性耦合例子，连同程序、输入、结果与收据，冻结于
\href{https://github.com/130U/certified-valuation-arithmetic-asian-options/tree/1959f8f078e1fedd590000d6166cb20429ff31db}{提交 \texttt{1959f8f078e1fedd590000d6166cb20429ff31db}}。
使用发表版本的检出目录及同一 Windows 环境，并指定新的目录名：

\begin{verbatim}
./.venv/Scripts/python.exe code/revision/round2/asian_parameter_point.py prepare --directory new-point
./.venv/Scripts/python.exe code/revision/round2/asian_parameter_point.py run --directory new-point
\end{verbatim}

运行在新目录中写入结果封包。
该参数点配置仅将 \(v_0\) 改为 \(1/25\)。
准备阶段检查原始数值清单，并构造显式参数副本。
加权、线性、试运行以及两个替代构造任务重新计算所有受影响量。
已声明的工作进程预算为：每项加权计算 120 秒，每项线性计算 360 秒，试运行 30 秒；一个线程和 256 MiB 的限制保持不变。
已记录的五任务外部总耗时约为 97.5 秒。
档案 \texttt{code/revision/round2/results/evidence-point-v004.zip} 包含全部 52 个执行成员及其清单。
其 SHA-256 和标准库复核说明见 \texttt{code/revision/round2/README.md}。
六项组成的精确宽度账本见 \texttt{code/revision/round2/results/point-v004-result.json}。

确定性耦合例子使用 384 位和 512 位精度的有限闭式高斯矩。
重新计算命令为

\begin{verbatim}
./.venv/Scripts/python.exe code/revision/round2_coupling_example.py --output coupling.json
\end{verbatim}

精确结果为 \texttt{code/revision/round2/results/coupling-v004-result.json}。
自含证明以及源码、结果哈希收据见 \texttt{docs/revision-round2/}。
独立的标量协方差矩计算也与这些包络一致。
该比较验证 \(\xi=0\) 上的非线性部分，不替代正 \(\xi\) 的定价计算。

\section{研究时间线}

研究与初稿写作于2023--2024年进行。
2026年完成投稿前的整理与核验。

\section{直接耦合余项不等式}

下述充分不等式通过耦合条件变量，保留两个条件余项之间的抵消。

\textbf{命题 I.1（耦合下的余项界）。}
固定 \(K>0\) 和 \(c>0\)，假设两个边际概率律均满足 \(\mathsf H_s(c)\)，从而余项有限。
耦合条件三元组 \((a_P,b_P,\sigma_P)\)、\((a_Q,b_Q,\sigma_Q)\)，其中 \(b_\nu=cg_\nu\)，所有分量均为正，且各边际均为式 (2.6) 中的条件三元组。
令
\[
s_- =\min(\sigma_P,\sigma_Q),\qquad
s_+ =\max(\sigma_P,\sigma_Q),\qquad
z=\max(|a_P-b_P|,|a_Q-b_Q|),
\]
\[
M_K=\frac{e^{2s_+^2}}{K s_-\sqrt{2\pi}},\qquad
N_K=\frac{e^{2s_+^2}(1+4/e+8s_+^2)}{K s_-^2\sqrt{2\pi}},
\]
\[
\Xi_K=M_Kz\left(|a_Q-a_P|+|b_Q-b_P|\right)
       +\frac12N_Kz^2|\sigma_Q-\sigma_P|.
\]
若 \(E\Xi_K\le B_K<\infty\)，则
\[
|R_{K,Q}-R_{K,P}|\le B_K.                       \tag{I.1}
\]
特别地，记 \(B=e^{-rT}(B_{K_1}+B_{K_2})\)，则看涨价差价格差中的精确非线性贡献属于
\([-B,B]\)。
在加入线性变换区间之前，可先将该区间与式 (4.2) 中的非线性区间取交集。
对相同条件概率律的对角耦合，有 \(\Xi_K=0\)。

\textit{证明。}
令 \(Z\) 为标准正态变量，并定义
\[
F_K(x;s)=E[(xe^{sZ}-K)^+],\qquad
r_K(a,b,s)=F_K(a;s)-F_K(b;s)-(a-b)\partial_xF_K(b;s).
\]
条件积分给出 \(R_{K,\nu}=E r_K(a_\nu,b_\nu,\sigma_\nu)\)。
对 \(x,s>0\)，记 \(y=\log(x/K)\)，则
\[
H_K(x,s):=\partial_x^2F_K(x;s)
=\frac{K}{s x^2\sqrt{2\pi}}e^{-y^2/(2s^2)},\qquad
\partial_sH_K=\frac{H_K}{s}\left(\frac{y^2}{s^2}-1\right).
\]
配方得到
\[
H_K(x,s)=\frac{e^{2s^2}}{K s\sqrt{2\pi}}
\exp\left[-\frac12\left(\frac ys+2s\right)^2\right].
\]
对 \(w=y/s+2s\)，有
\[
\left|(w-2s)^2-1\right|e^{-w^2/2}
\le(2w^2+8s^2+1)e^{-w^2/2}
\le 4/e+8s^2+1.
\]
因此，只要 \(s\in[s_-,s_+]\)，就有
\(0\le H_K\le M_K\) 和 \(|\partial_sH_K|\le N_K\)，且关于
\(x>0\) 一致成立。
对余项求导并使用 Taylor 积分公式，得到
\[
\partial_a r_K=\partial_xF_K(a;s)-\partial_xF_K(b;s),\qquad
\partial_b r_K=-(a-b)H_K(b,s),
\]
\[
\partial_s r_K=(a-b)^2\int_0^1(1-t)
       \partial_sH_K(b+t(a-b),s)\,dt.
\]
在两个耦合端点之间，对三个分量作线性插值。
所有插值后的 \(a,b,s\) 仍为正，且
\(|a-b|\le z\)、\(s\in[s_-,s_+]\)。
由微积分基本定理及上述导数界，对每个耦合对均有
\[
|r_K(a_Q,b_Q,\sigma_Q)-r_K(a_P,b_P,\sigma_P)|\le\Xi_K.
\]
取期望即证明式 (I.1)。
对两个执行价应用三角不等式并乘以折现因子，即得以价格为单位的结论。
\hfill\(\square\)

因子 \(s_-^{-2}\) 需要超出式 (2.7) 和边际仿射项列表的耦合矩控制。
若一族经过验证的耦合满足 \(E\Xi_K\to0\)，则式 (I.1) 给出趋零的余项差。
本文不对正波动率之波动率的 Heston 例子宣称这样的收敛界或收敛速度。

\subsection{确定性方差参数面上已执行的耦合例子}

在 \(\xi=0,\kappa=3,\bar v=9/200,v_0=1/25,S_0=100,r=1/100,T=1,h=1/768\) 下，保留十二个按月观察时点、执行价为 95 和 110 的看涨价差，以及
\(c=1254433/1250000\)。
记
\(s_{\nu,j}=I_\nu(t_j)-I_\nu(t_{j-1})\)，其中使用式 (6.2) 中的精确连续及 Euler 累计方差。
通过十二个共同独立标准正态变量耦合两种概率律。
令 \(\sigma_\nu=\sqrt{s_{\nu,1}}\)，分离第一个中心化正态因子，并置
\[
m_{\nu,i}=rt_i-I_\nu(t_i)/2,\qquad
\ell_{\nu,i,j}=\sqrt{s_{\nu,j}}\mathbf1_{\{j\le i\}},\quad 2\le j\le12,
\]
\[
a_{\nu,i}=100e^{m_{\nu,i}+\ell_{\nu,i}\cdot Z},\qquad
a_\nu=\frac1{12}\sum_i a_{\nu,i},\qquad
g_\nu=100e^{\bar m_\nu+\bar\ell_\nu\cdot Z},\qquad b_\nu=cg_\nu.
\]
这里，\(Z=(Z_2,\ldots,Z_{12})\)，上横线表示对观察时点指标 \(i=1,\ldots,12\) 取平均。
漂移包含首月项 \(-s_{\nu,1}/2\)，因此分离出的因子具有 \(\mathsf H_s(c)\) 所要求的中心化高斯分布。
这里将股价驱动 \(\rho W+\sqrt{1-\rho^2}B\) 合成为一个布朗运动。
在 \(\xi=0\) 上，方差是确定性的，价格概率律与 \(\rho\) 无关；因此 \(\sigma_\nu^2=s_{\nu,1}\)，使用首月完整方差。
下面的分模型比较使用相同的条件化。

所有条件二阶矩，包括混合 P/Q 矩，均由下式给出：
\[
E\!\left[100e^{m+u\cdot Z}\,100e^{n+v\cdot Z}\right]
=10000e^{m+n+\|u+v\|^2/2}.
\]
定义
\[
q_\nu=E(a_\nu-b_\nu)^2,\quad Z_2=q_P+q_Q,\quad
D_a=E(a_Q-a_P)^2,\quad D_b=E(b_Q-b_P)^2.
\]
对耦合命题中确定性的 \(M_K,N_K\)，由逐点界 \(z^2\le(a_P-b_P)^2+(a_Q-b_Q)^2\) 和 Cauchy--Schwarz 不等式，得到可执行上界
\[
E\Xi_K\le B_K:=
M_K\sqrt{Z_2}(\sqrt{D_a}+\sqrt{D_b})
+\frac12N_KZ_2|\sigma_Q-\sigma_P|.
\]
384 位与 512 位精度的有限高斯矩计算给出
\[
E\Xi_{95}\le .000023887739528734380066,\qquad
E\Xi_{110}\le .000020630320502088782784.
\]
因此，折现后的非线性价差贡献属于半径为 \(e^{-r}(B_{95}+B_{110})\) 的对称区间，其宽度至多为 \(.000088150195864703911227\)。
对相同的概率律、条件化和缩放系数，分别控制两模型的估计给出的宽度位于
\[
[.015383637151510699801154,\;.015383637151510699801155].
\]
两宽度之比低于 \(.005731\)。
严格改进通过耦合宽度的向外上端点与分别估计宽度的下端点进行检查。
比较使用
\(D_\nu=e^{2\sigma_\nu^2}q_\nu/\sigma_\nu\)，因此所有加权矩都对应同一高斯因子。

这是通过闭式矩给出的耦合余项期望包络，不使用蒙特卡洛或无界域求积。
它证明了单个确定性方差参数点上非线性部分的收紧。
它不计算完整价格区间、有符号的亚式首阶系数，或 \(\xi>0\) 的耦合证书。
两种精度和两种高斯矩代数表达式共享解析证明与 Arb；它们是算术一致性检查。

\begin{thebibliography}{13}
\bibitem{Curran1994}
Curran, Michael.
Valuing {Asian} and Portfolio Options by Conditioning on the Geometric Mean Price.
\emph{Management Science}, 40(12):1705--1711, 1994.
\url{https://doi.org/10.1287/mnsc.40.12.1705}.

\bibitem{RogersShi1995}
Rogers, L. C. G., Shi, Z..
The value of an {Asian} option.
\emph{Journal of Applied Probability}, 32(4):1077--1088, 1995.
\url{https://doi.org/10.2307/3215221}.

\bibitem{Thompson2002}
Thompson, G. W. P..
Fast narrow bounds on the value of {Asian} options.
Judge Institute of Management, University of Cambridge, WP09/2002, 2002.
\url{https://www.jbs.cam.ac.uk/wp-content/uploads/2020/08/wp0209.pdf}.

\bibitem{FusaiKyriakou2016}
Fusai, Gianluca, and Kyriakou, Ioannis.
General Optimized Lower and Upper Bounds for Discrete and Continuous Arithmetic {Asian} Options.
\emph{Mathematics of Operations Research}, 41(2):531--559, 2016.
\url{https://doi.org/10.1287/moor.2015.0739}.
Accepted manuscript: \url{https://openaccess.city.ac.uk/id/eprint/13241/1/AsianBound_FK.pdf}.

\bibitem{Heston1993}
Heston, Steven L..
A Closed-Form Solution for Options with Stochastic Volatility with Applications to Bond and Currency Options.
\emph{The Review of Financial Studies}, 6(2):327--343, 1993.
\url{https://doi.org/10.1093/rfs/6.2.327}.

\bibitem{KimKimKimWee2016}
Kim, Bara, Kim, Jeongsim, Kim, Jerim, and Wee, In Suk.
A recursive method for discretely monitored geometric {Asian} option prices.
\emph{Bulletin of the Korean Mathematical Society}, 53(3):733--749, 2016.
\url{https://doi.org/10.4134/BKMS.b150283}.

\bibitem{QuantLibDiscreteGeometricHeston}
{QuantLib contributors}.
{AnalyticDiscreteGeometricAveragePriceAsianHestonEngine}.
Official source code; fixed revision.
\url{https://github.com/lballabio/QuantLib/blob/966a4cc101049ca36a888b2ce223aa96d3f3b22d/ql/experimental/asian/analytic_discr_geom_av_price_heston.hpp}.

\bibitem{TalayTubaro1990}
Talay, Denis, Tubaro, Luciano.
Expansion of the global error for numerical schemes solving stochastic differential equations.
\emph{Stochastic Analysis and Applications}, 8(4):483--509, 1990.
\url{https://doi.org/10.1080/07362999008809220}.

\bibitem{MickelNeuenkirch2023}
Mickel, Annalena, and Neuenkirch, Andreas.
The weak convergence order of two {Euler}-type discretization schemes for the log-{Heston} model.
\emph{IMA Journal of Numerical Analysis}, 43(6):3326--3356, 2023.
\url{https://doi.org/10.1093/imanum/drac069}.

\bibitem{BallyTalay1996}
Bally, Vlad, Talay, Denis.
The law of the {Euler} scheme for stochastic differential equations: {I}. {Convergence} rate of the distribution function.
\emph{Probability Theory and Related Fields}, 104(1):43--60, 1996.
\url{https://doi.org/10.1007/BF01303802}.

\bibitem{CotterDashtiStuart2010}
Cotter, S. L., Dashti, M., Stuart, A. M..
Approximation of {Bayesian} Inverse Problems for {PDEs}.
\emph{SIAM Journal on Numerical Analysis}, 48(1):322--345, 2010.
\url{https://doi.org/10.1137/090770734}.

\bibitem{Johansson2017}
Johansson, Fredrik.
{Arb}: Efficient Arbitrary-Precision Midpoint-Radius Interval Arithmetic.
\emph{IEEE Transactions on Computers}, 66(8):1281--1292, 2017.
\url{https://doi.org/10.1109/TC.2017.2690633}.

\bibitem{Lee2004}
Lee, Roger W..
Option Pricing by Transform Methods: Extensions, Unification, and Error Control.
\emph{The Journal of Computational Finance}, 7(3):51--86, 2004.
\url{https://doi.org/10.21314/JCF.2004.121}.

\end{thebibliography}
